\HeaderA{Details}{Lagrangian Multiplier Smoothing Splines: Mathematical Details}{Details}
%
\begin{Description}
This document provides the mathematical and implementation details for
Lagrangian Multiplier Smoothing Splines as implemented in \pkg{lgspline}.

This package provides exhaustive resources for fitting multivariate smoothing
smoothing splines with a monomial basis and analytical form cubic smoothing
spline penalty.

The material is presented such that a programmer or statistician of
reasonable experience and background can understand and implement the
procedure from scratch, and also potentially critique some of the modelling
choices that went into designing this package.

Informally, \pkg{lgspline} answers the following question:
How can we best adapt a useful functionality of basis splines, under the
alternative interpretation of smoothing as explicit external constraints instead?

The obvious benefit is a much more flexible and interpretable final model that for
non-experienced users is simply easier to understand without post-hoc processing, and
for experienced users can be used to customize models more easily.

The drawback is that the interpretation of constraints as external adds a new layer of complexity
to each step of the model fitting process, whereas for implicit design matrix
construction these complications are bypassed.

While it is true a B-spline can always be converted back into monomial form,
tensor-product splines that generalize this to multiple dimensions often
explodes the number and degree of interaction terms, the conversion may not
be computationally stable, and it is not available in standard software.
\end{Description}
%
\begin{Section}{Statistical Problem Formulation}


Consider an \eqn{N \times q}{} matrix of predictors
\eqn{\mathbf{T} = (\mathbf{t}_1, \dots, \mathbf{t}_N)^{\top}}{} and an
\eqn{N \times 1}{} response vector \eqn{\mathbf{y} = (y_1, \dots, y_N)^{\top}}{}.
We assume the relationship follows a generalized linear model with unknown
smooth function \eqn{f}{}:
\deqn{y_i \sim \mathcal{D}(g^{-1}(f(\mathbf{t}_i)),\, \sigma^2)}{}
where \eqn{\mathcal{D}}{} is a distribution (e.g. exponential family or related) with mean
\eqn{\mu_i = g^{-1}(f(\mathbf{t}_i))}{}, link function \eqn{g(\cdot)}{}, and
dispersion parameter \eqn{\sigma^2}{}. For Gaussian response with identity
link, observations are independently distributed as
\eqn{y_i \mid \mathbf{t}_i, \sigma^2 \sim \mathcal{N}(f(\mathbf{t}_i), \sigma^2)}{}.

The objective is to estimate \eqn{f}{} by:
\begin{enumerate}

\item{} Partitioning the predictor space into \eqn{K+1}{} mutually exclusive regions.
\item{} Fitting local polynomial models within each partition.
\item{} Enforcing smoothness at partition boundaries via Lagrangian multipliers.
\item{} Penalizing the integrated squared second derivative to discourage roughness.

\end{enumerate}


Unlike other smoothing spline formulations, no post-fitting algebraic
rearrangement or disentanglement of a spline basis is needed to obtain
interpretable models. The polynomial expansions are homogeneous across
partitions, and the relationship between predictor and response is explicit
at the coefficient level.

To anchor the notation, in the single-predictor cubic case one would write
\deqn{\hat{f}(t_i) = \hat{\beta}_{(0)} + \hat{\beta}_{(1)}t_i +
  \hat{\beta}_{(2)}t_i^2 + \hat{\beta}_{(3)}t_i^3 =
  \mathbf{x}_i^{\top}\hat{\boldsymbol{\beta}},}{}
where \eqn{\mathbf{x}_i = (1, t_i, t_i^2, t_i^3)^{\top}}{}. The LMSS
formulation preserves exactly this kind of polynomial representation, but now
does so within each partition and then forces neighboring pieces to agree in
the smoothness conditions described below.

Core notation used throughout:
\begin{itemize}

\item{} \eqn{\mathbf{y}_{(N \times 1)}}{}: Response vector.
\item{} \eqn{\mathbf{T}_{(N \times q)}}{}: Matrix of predictors.
\item{} \eqn{\mathbf{X}_{(N \times P)}}{}: Block-diagonal matrix of polynomial
expansions, with diagonal blocks \eqn{\mathbf{X}_k}{} of dimension
\eqn{n_k \times p}{}.
\item{} \eqn{\boldsymbol{\Lambda}_{(P \times P)}}{}: Block-diagonal penalty matrix,
with blocks \eqn{\boldsymbol{\Lambda}_k}{} of dimension \eqn{p \times p}{}.
\item{} \eqn{\hat{\boldsymbol{\beta}}_{(P \times 1)}}{}: Unconstrained penalized
estimate.
\item{} \eqn{\tilde{\boldsymbol{\beta}}_{(P \times 1)}}{}: Constrained coefficient
estimates.
\item{} \eqn{\mathbf{G}_{(P \times P)}}{}: Block-diagonal matrix with blocks
\eqn{\mathbf{G}_k = (\mathbf{X}_k^{\top}\mathbf{W}_k\mathbf{D}_k\mathbf{X}_k + \boldsymbol{\Lambda}_k)^{-1}}{},
where \eqn{\mathbf{W}_k}{} and \eqn{\mathbf{D}_k}{} are defined below.
\item{} \eqn{\mathbf{A}_{(P \times r)}}{}: Constraint matrix encoding smoothness
conditions. Reduced to linearly independent columns via pivoted QR
decomposition.
\item{} \eqn{\mathbf{U}_{(P \times P)}}{}:
\eqn{\mathbf{I} - \mathbf{G}\mathbf{A}(\mathbf{A}^{\top}\mathbf{G}\mathbf{A})^{-1}\mathbf{A}^{\top}}{}.
\item{} \eqn{\mathbf{D}_{(N \times N)}}{}: Diagonal matrix of user-supplied
observation weights (\code{observation\_weights} or \code{weights}).
Defaults to the identity. These play the role of prior precision on
individual observations: a weight of 2 is equivalent to seeing that
observation twice.
\item{} \eqn{\mathbf{W}_{(N \times N)}}{}: Diagonal matrix of GLM working
weights. In the implementation these diagonal entries are whatever is
returned by \code{glm\_weight\_function}; by default this is
\code{family\$variance(mu)}, optionally multiplied by user-supplied
observation weights. For Gaussian response with identity link,
\eqn{\mathbf{W} = \mathbf{I}}{}. For other families,
\eqn{\mathbf{W}}{} depends on the current fitted values and is updated at
each Newton-Raphson iteration. For the common canonical families used
by default, this matches the familiar Fisher-scoring weighting role.
\item{} \eqn{\mathbf{V}_{(N \times N)}}{}: Correlation matrix of errors.
When no correlation structure is specified, \eqn{\mathbf{V} = \mathbf{I}}{}.
Otherwise supplied via \code{VhalfInv} or estimated through
\code{VhalfInv\_fxn}.

\end{itemize}


In the Gaussian identity case with unit weights and no correlation,
\eqn{\mathbf{G}_k = (\mathbf{X}_k^{\top}\mathbf{X}_k + \boldsymbol{\Lambda}_k)^{-1}}{}
and most formulas simplify accordingly. When \eqn{\mathbf{D}}{} or
\eqn{\mathbf{W}}{} appear in a formula, the product \eqn{\mathbf{W}\mathbf{D}}{}
means ``GLM working weights times observation weights''; whenever one of
them is the identity it drops out.

Before these quantities reach the main fitting stage, the user-facing inputs
are parsed, standardized, and organized by \code{\LinkA{process\_input}{process.Rul.input}}. When
the formula interface is used and \code{auto\_encode\_factors = TRUE}, that
preprocessing step also relies on helpers such as \code{\LinkA{create\_onehot}{create.Rul.onehot}}
to encode factor levels before the design reaches \code{\LinkA{lgspline.fit}{lgspline.fit}()}.
The notation in the remainder of this document therefore refers to the internal
objects that actually enter \code{\LinkA{lgspline.fit}{lgspline.fit}()}, not necessarily the raw
objects originally supplied by the user.

From the user side, many of the arguments that control these internal objects
can be supplied either individually or through the grouped lists
\code{penalty\_args}, \code{tuning\_args}, \code{expansion\_args},
\code{constraint\_args}, \code{qp\_args}, \code{parallel\_args},
\code{covariance\_args}, \code{return\_args}, and \code{glm\_args}, as
documented in \code{\LinkA{lgspline}{lgspline}}. These grouped lists are unpacked before
dispatch into the same fitting pipeline, so they are a convenience layer
rather than a separate modeling abstraction. A closely related exploratory
mode is \code{dummy\_fit = TRUE} in \code{\LinkA{lgspline}{lgspline}} or
\code{\LinkA{lgspline.fit}{lgspline.fit}}, which runs the preprocessing, partition
construction, expansion building, and penalty setup without solving for
nonzero coefficients, making it a practical way to inspect objects such as
\code{X}, \code{A}, the returned \code{make\_partition\_list} from
\code{\LinkA{make\_partitions}{make.Rul.partitions}}, and the assembled \code{penalties} from
\code{\LinkA{compute\_Lambda}{compute.Rul.Lambda}} before a full fit.
\end{Section}
%
\begin{Section}{Model Formulation and Estimation}


%
\begin{SubSection}{Piecewise Polynomial Structure}
For \eqn{K}{} knots (one predictor) or \eqn{K+1}{} partitions (multiple
predictors) there are \eqn{K+1}{} mutually exclusive partitions
\eqn{\mathcal{P}_0, \dots, \mathcal{P}_{K}}{}. Each observation \eqn{i}{}
belongs to exactly one partition. Within partition \eqn{k}{}, the function is
represented as a polynomial of degree \eqn{p-1}{} in each predictor:
\deqn{\hat{f}_k(\mathbf{t}) = \mathbf{x}^{\top}\tilde{\boldsymbol{\beta}}_k}{}
where \eqn{\mathbf{x}}{} collects the polynomial basis terms (intercept,
linear, quadratic, cubic, and optionally quartic and interaction terms) and
\eqn{\tilde{\boldsymbol{\beta}}_k}{} are the corresponding coefficients. In
one predictor, the same idea can be written more explicitly as
\deqn{\hat{f}(t_i) = \sum_{k=0}^{K}
  \mathbf{x}_{ik}^{\top}\hat{\boldsymbol{\beta}}_k
  \mathbf{1}(t_i \in \mathcal{P}_k),}{}
which highlights that the unconstrained problem is just a collection of
local polynomial regressions. The expansions are homogeneous across
partitions, so coefficients are directly comparable. This is implemented via
\code{\LinkA{get\_polynomial\_expansions}{get.Rul.polynomial.Rul.expansions}}.

The exact contents of \eqn{\mathbf{x}}{} are controlled by the basis-expansion
arguments documented in \code{\LinkA{lgspline}{lgspline}}: \code{include\_quadratic\_terms},
\code{include\_cubic\_terms}, \code{include\_quartic\_terms},
\code{include\_2way\_interactions}, \code{include\_3way\_interactions},
\code{include\_quadratic\_interactions}, \code{exclude\_interactions\_for},
\code{exclude\_these\_expansions}, and \code{custom\_basis\_fxn}. Likewise,
\code{just\_linear\_with\_interactions} and
\code{just\_linear\_without\_interactions} determine which predictors remain
structurally linear even though they still participate in the same
partition-wise polynomial bookkeeping described here.

Letting \eqn{p}{} denote the number of basis terms per partition,
\eqn{P = p(K+1)}{} is the total number of coefficients. The full design matrix
\eqn{\mathbf{X}}{} and penalty matrix \eqn{\boldsymbol{\Lambda}}{} are
block-diagonal with \eqn{K+1}{} blocks, so unconstrained estimation reduces to
\eqn{K+1}{} independent penalized regressions, which appears as follows for the identity link case:
\deqn{\hat{\boldsymbol{\beta}}_k = \mathbf{G}_k \mathbf{X}_k^{\top} \mathbf{W}_k\mathbf{D}_k\mathbf{y}_k, \quad
  \mathbf{G}_k = (\mathbf{X}_k^{\top}\mathbf{W}_k\mathbf{D}_k\mathbf{X}_k + \boldsymbol{\Lambda}_k)^{-1}.}{}
For Gaussian identity with unit weights this reduces to the familiar
\eqn{\mathbf{G}_k = (\mathbf{X}_k^{\top}\mathbf{X}_k + \boldsymbol{\Lambda}_k)^{-1}}{}.
The block-diagonal structure means these can be computed in parallel across
partitions. In the user-facing interface this is realized by supplying a
cluster through \code{cl}, optionally controlling work splitting with
\code{chunk\_size}, and enabling stages such as \code{parallel\_eigen}
for the eigendecompositions and, in non-Gaussian Path 3, \code{parallel\_unconstrained}
for the partition-wise unconstrained fits; nearby stages can likewise use
\code{parallel\_penalty} and \code{parallel\_make\_constraint}. The eigenvalue
decomposition and matrix square roots of each \eqn{\mathbf{G}_k}{} are
computed by \code{\LinkA{compute\_G\_eigen}{compute.Rul.G.Rul.eigen}}, and can be returned in the fitted
object as \code{G} and \code{Ghalf} when \code{return\_G = TRUE} and
\code{return\_Ghalf = TRUE}.

Fitted values for the canonical Gaussian case appear as
\eqn{\tilde{\mathbf{y}} = \mathbf{X}\tilde{\boldsymbol{\beta}} = \mathbf{H}\mathbf{y}}{}
for \eqn{\mathbf{H} = \mathbf{X}\mathbf{U}\mathbf{G}\mathbf{X}^{\top}}{}.
\end{SubSection}


%
\begin{SubSection}{Smoothing Constraints and the Constraint Matrix}
Without further intervention the piecewise polynomial will be discontinuous.
The central idea of LMSS is that smoothness is not hidden inside a special
basis, but instead imposed directly where neighboring partitions meet.
At each knot \eqn{t_{k,k+1}}{} between neighboring partitions \eqn{k}{} and
\eqn{k+1}{}, up to three smoothing constraints are imposed:
\begin{enumerate}

\item{} Continuity: \eqn{\mathbf{x}_{k,k+1}^{\top}\boldsymbol{\beta}_k = \mathbf{x}_{k,k+1}^{\top}\boldsymbol{\beta}_{k+1}}{}.
\item{} First-derivative continuity: \eqn{\mathbf{x}_{k,k+1}^{\prime\top}\boldsymbol{\beta}_k = \mathbf{x}_{k,k+1}^{\prime\top}\boldsymbol{\beta}_{k+1}}{}.
\item{} Second-derivative continuity: \eqn{\mathbf{x}_{k,k+1}^{\prime\prime\top}\boldsymbol{\beta}_k = \mathbf{x}_{k,k+1}^{\prime\prime\top}\boldsymbol{\beta}_{k+1}}{}.

\end{enumerate}

where \eqn{\mathbf{x}^{\prime}}{} and \eqn{\mathbf{x}^{\prime\prime}}{} are
elementwise first and second derivatives of the basis with respect to
\eqn{\mathbf{t}}{}. For the familiar cubic single-predictor basis
\eqn{\mathbf{x} = (1, t, t^2, t^3)^{\top}}{}, these derivative vectors are
\deqn{\mathbf{x}' = (0, 1, 2t, 3t^2)^{\top}, \qquad
  \mathbf{x}'' = (0, 0, 2, 6t)^{\top}.}{}
With \eqn{K}{} knots this yields up to \eqn{3K}{} scalar
constraints (for a single predictor; more for multiple predictors with
interactions), collected as linear equations
\eqn{\mathbf{A}^{\top}\boldsymbol{\beta} = \mathbf{0}}{}
in a \eqn{P \times r}{} matrix \eqn{\mathbf{A}}{}. The constraint matrix is
built by \code{\LinkA{make\_constraint\_matrix}{make.Rul.constraint.Rul.matrix}} and returned in the fitted
object as \code{A}.

In higher dimensions or with many partitions, the constraints can become
over-specified and force the model toward a single global polynomial. In these
cases it is recommended to drop second-derivative constraints or include quartic
terms, allowing the model to fit a richer surface while maintaining
perceived smoothness at knots. The appropriate constraint level can be
controlled via \code{include\_constrain\_fitted},
\code{include\_constrain\_first\_deriv}, and
\code{include\_constrain\_second\_deriv}.
The companion flag \code{include\_constrain\_interactions} determines whether
the analogous mixed-partial constraints are imposed for interaction terms,
and \code{no\_intercept} adds the special homogeneous equality constraint that
fixes the intercept at zero (the same behavior triggered by using
\code{0 +} in the formula interface).

Before computing the projection \eqn{\mathbf{U}}{}, the constraint matrix is
reduced to a linearly independent subset of columns via pivoted QR
decomposition. This avoids numerical instability from redundant constraints
and ensures \eqn{\mathbf{A}^{\top}\mathbf{G}\mathbf{A}}{} is invertible.
\end{SubSection}


%
\begin{SubSection}{Lagrangian Projection}
The constrained estimate is derived via Lagrangian multipliers. Define the
\eqn{P \times P}{} projection matrix:
\deqn{\mathbf{U} = \mathbf{I} - \mathbf{G}\mathbf{A}(\mathbf{A}^{\top}\mathbf{G}\mathbf{A})^{-1}\mathbf{A}^{\top}.}{}
Then the constrained estimate is:
\deqn{\tilde{\boldsymbol{\beta}} = \mathbf{U}\hat{\boldsymbol{\beta}}.}{}
The matrix \eqn{\mathbf{U}}{} has the property that
\eqn{\mathbf{U}\mathbf{G}\mathbf{U}^{\top} = \mathbf{U}\mathbf{G}}{}, which is
used extensively in variance estimation and posterior draws. In words, the
unconstrained penalized estimate is projected back into the coefficient space
that satisfies the smoothness restrictions, and all subsequent uncertainty
calculations inherit that same projected geometry.
The projection is computed via \code{\LinkA{get\_U}{get.Rul.U}} and, when requested,
returned in the fitted object as \code{U} through \code{return\_U = TRUE}.

When the constraints are inhomogeneous
(\eqn{\mathbf{A}^{\top}\boldsymbol{\beta} = \mathbf{c}}{} with
\eqn{\mathbf{c} \neq \mathbf{0}}{}), a particular solution
\eqn{\boldsymbol{\beta}_0}{} satisfying
\eqn{\mathbf{A}^{\top}\boldsymbol{\beta}_0 = \mathbf{c}}{} is added back
after projection, yielding the full Lagrangian solution
\eqn{\mathbf{U}\hat{\boldsymbol{\beta}} + (\mathbf{I} - \mathbf{U})\boldsymbol{\beta}_0}{}.
In \code{\LinkA{lgspline}{lgspline}} and \code{\LinkA{lgspline.fit}{lgspline.fit}}, users realize
this by supplying extra equality columns in \code{constraint\_vectors}
together with matching right-hand sides in \code{constraint\_values};
\code{null\_constraint} provides the alternate shorthand documented in
\code{\LinkA{lgspline}{lgspline}} when \code{constraint\_vectors} is supplied and
\code{constraint\_values} is left empty.

In practice \eqn{\mathbf{U}}{} is never explicitly formed during fitting.
The constrained estimate is obtained from a transformed OLS residual problem (the
\eqn{\mathbf{G}^{1/2}\mathbf{r}^{*}}{} trick) in four steps:
\begin{enumerate}

\item{} Obtain the unconstrained partition-wise
unconstrained estimate \eqn{\hat{\boldsymbol{\beta}}}{}.
\item{} Set \eqn{\mathbf{y}^{*} = \mathbf{G}^{-1/2}\hat{\boldsymbol{\beta}}}{} and
\eqn{\mathbf{X}^{*} = \mathbf{G}^{1/2}\mathbf{A}}{}.
\item{} Fit the linear model \eqn{\mathbb{E}[\mathbf{y}^{*}] = \mathbf{X}^{*}\boldsymbol{\gamma}}{}
by OLS using QR decomposition.
\item{} Compute the residuals \eqn{\mathbf{r}^{*} = \mathbf{y}^{*} - \mathbf{X}^{*}(\mathbf{X}^{*\top}\mathbf{X}^{*})^{-1}\mathbf{X}^{*\top}\mathbf{y}^{*}}{} from that transformed OLS
fit and recover the constrained estimate by
\eqn{\tilde{\boldsymbol{\beta}} = \mathbf{G}^{1/2}\mathbf{r}^{*}}{}.

\end{enumerate}


A scaling factor \eqn{1/\sqrt{K+1}}{} is applied to both
\eqn{\mathbf{X}^{*}}{} and \eqn{\mathbf{y}^{*}}{} prior to the OLS call
and divided out afterward, improving numerical conditioning when the
constraint matrix has many rows.

The most expensive operation in this approach is the QR decomposition of the
\eqn{P \times r}{} matrix \eqn{\mathbf{X}^{*} = \mathbf{G}^{1/2}\mathbf{A}}{},
which is far cheaper than working with the full \eqn{P \times P}{} system directly.
Without correlation or SQP constraints, \eqn{\mathbf{G}}{} is stored and operated upon as a
list of \eqn{K+1}{} small \eqn{p \times p}{} matrices rather than the full
\eqn{P \times P}{} block-diagonal, saving substantial memory when \eqn{K}{}
is large and allowing for parallelism.

When correlation is present, \eqn{\mathbf{X}^{\top}\mathbf{V}^{-1}\mathbf{X}}{}
is no longer block-diagonal, so the full-dimensional system must be handled
directly unless the Woodbury acceleration
(see \code{\LinkA{.woodbury\_decompose\_V}{.woodbury.Rul.decompose.Rul.V}}) applies.
When additional inequality constraints are present, the code either augments the equality system with a
partition-wise active-set refinement (block-separable case) or falls back to
dense SQP via the Goldfarb-Idnani dual active-set method implemented in
\code{\LinkA{solve.QP}{solve.QP}}.
\end{SubSection}

\end{Section}
%
\begin{Section}{GLM Extension and Iterative Updates}


%
\begin{SubSection}{Working Quantities}
For GLM responses with mean \eqn{\mu_i = g^{-1}(\eta_i)}{} and working
weights \eqn{w_i = [V(\mu_i)\{g'(\mu_i)\}^2]^{-1}}{}, the penalized
information becomes
\eqn{\mathbf{G}_k^{-1} = \mathbf{X}_k^{\top}\mathbf{W}_k\mathbf{X}_k
+ \boldsymbol{\Lambda}_k}{} with
\eqn{\mathbf{W}_k = \mathrm{diag}(w_i : i \in \mathcal{P}_k)}{}.
Because \eqn{\mathbf{W}}{} is diagonal,
\eqn{\mathbf{X}^{\top}\mathbf{W}\mathbf{X}}{} remains block-diagonal and
the four-step procedure carries through with
\eqn{\mathbf{G}_k = (\mathbf{X}_k^{\top}\mathbf{W}_k\mathbf{X}_k +
\boldsymbol{\Lambda}_k)^{-1}}{} in place of
\eqn{(\mathbf{X}_k^{\top}\mathbf{X}_k + \boldsymbol{\Lambda}_k)^{-1}}{}.
Under the default \code{glm\_weight\_function}, the diagonal entries reduce
to \code{family\$variance(mu)} (optionally scaled by observation weights),
which for canonical families is the usual Fisher scoring weight.

The partition-wise unconstrained estimates are obtained by
\code{unconstrained\_fit\_fxn}, by default
\code{\LinkA{unconstrained\_fit\_default}{unconstrained.Rul.fit.Rul.default}}, which initializes via an augmented
ridge trick (appending \eqn{\boldsymbol{\Lambda}^{1/2}}{} as
pseudo-observations to \code{glm.fit}) then refines using
\code{\LinkA{damped\_newton\_r}{damped.Rul.newton.Rul.r}} with \code{\LinkA{nr\_iterate}{nr.Rul.iterate}}. For
non-canonical log-link gamma regression, the \code{keep\_weighted\_Lambda = TRUE}
option correctly returns the augmented ridge estimate directly.
\end{SubSection}


%
\begin{SubSection}{Three Fitting Paths}

\strong{Path 1: Correlation structure present.}
When a working correlation \eqn{\mathbf{V}}{} is present, such as for
marginal means models or generalized estimating equation (GEE)-like models,
\eqn{\mathbf{V}^{-1}}{} couples observations across partitions, so
partition-wise fitting is not directly available. Two sub-paths handle
this.

\emph{Path 1a} (Gaussian identity + GEE) solves the whitened system
directly via
\eqn{\tilde{\mathbf{G}} =
(\tilde{\mathbf{X}}^{\top}\tilde{\mathbf{X}} +
\boldsymbol{\Lambda}_{\mathrm{block}})^{-1}}{} where
\eqn{\tilde{\mathbf{X}} = \mathbf{V}^{-1/2}\mathbf{X}}{}, then applies
the Lagrangian projection in the full \eqn{P}{}-space.
See \code{\LinkA{.get\_B\_gee\_gaussian}{.get.Rul.B.Rul.gee.Rul.gaussian}}.

\emph{Path 1b} (non-Gaussian GEE) uses a damped SQP iteration on the
whitened system. The first iterate is a constrained Newton step from the
projection matrix \eqn{\mathbf{U}}{}; subsequent iterates solve the
quadratic subproblem via \code{\LinkA{solve.QP}{solve.QP}}.
See \code{\LinkA{.get\_B\_gee\_glm}{.get.Rul.B.Rul.gee.Rul.glm}}.

Both sub-paths have Woodbury-accelerated variants described below.

\strong{Path 2: Gaussian identity, no correlation.}
The constrained estimate is obtained by a single Lagrangian projection
from the per-partition penalized least-squares cross-products (the four
steps in closed form). No outer iteration is needed. When \eqn{K = 0}{} and
there are no additional constraints, this reduces to the ordinary
penalized closed form
\eqn{\hat{\boldsymbol{\beta}} = \mathbf{G}\mathbf{X}^{\top}\mathbf{y}}{}.
Implemented in \code{\LinkA{.get\_B\_gaussian\_nocorr}{.get.Rul.B.Rul.gaussian.Rul.nocorr}}.

\strong{Path 3: Non-Gaussian GLM, no correlation.}
Unconstrained estimates are obtained separately within each partition,
then projected onto the constraint space. When the link is non-identity,
the information \eqn{\mathbf{G}}{} depends on the current fitted values
through the working weights, so the projection must be iterated. The
unconstrained anchor \eqn{\hat{\boldsymbol{\beta}}}{} is held fixed while
\eqn{\mathbf{G}}{} is recomputed at each iterate's constrained estimate:
\deqn{\tilde{\boldsymbol{\beta}}^{(s+1)} =
  \mathbf{U}^{(s)}\hat{\boldsymbol{\beta}}, \quad
  \mathbf{U}^{(s)} = \mathbf{I} - \mathbf{G}^{(s)}\mathbf{A}
  (\mathbf{A}^{\top}\mathbf{G}^{(s)}\mathbf{A})^{-1}\mathbf{A}^{\top}.}{}
Iteration stops when the mean absolute coefficient change falls below
\code{tol}, or when the update begins increasing (in which case the
previous iterate is restored). The recomputation of weighted Gram
matrices, Schur corrections, and square-root information factors at each
step is handled by \code{.solver\_recompute\_G\_at\_estimate}. Implemented
in \code{\LinkA{.get\_B\_glm\_nocorr}{.get.Rul.B.Rul.glm.Rul.nocorr}}.
\end{SubSection}


%
\begin{SubSection}{Woodbury Acceleration for Structured Correlation}
For structured correlation, write the penalized information as
\deqn{\mathbf{G}^{-1}
  = \mathbf{G}_{\mathrm{on}}^{-1} + \mathbf{G}_{\mathrm{off}}^{-1},}{}
where \eqn{\mathbf{G}_{\mathrm{on}}^{-1}}{} contains the \eqn{K+1}{}
block-diagonal \eqn{p \times p}{} pieces of
\eqn{\mathbf{X}^{\top}\mathbf{V}^{-1}\mathbf{X} + \boldsymbol{\Lambda}}{},
and \eqn{\mathbf{G}_{\mathrm{off}}^{-1}}{} contains the cross-partition part.
Without correlation, \eqn{\mathbf{G}_{\mathrm{off}}^{-1} = \mathbf{0}}{}.

When the cross-partition part is low rank, the code factors it as
\deqn{\mathbf{G}_{\mathrm{off}}^{-1} = \mathbf{E}\mathbf{J}\mathbf{E}^{\top},}{}
with \eqn{\mathbf{E}}{} of size \eqn{P \times r}{} and \eqn{\mathbf{J}}{}
diagonal with entries \eqn{\pm 1}{}. Then
\deqn{\mathbf{G}
  = \mathbf{G}_{\mathrm{on}}
  - \mathbf{G}_{\mathrm{on}}\mathbf{E}
  (\mathbf{J}^{-1} + \mathbf{E}^{\top}\mathbf{G}_{\mathrm{on}}\mathbf{E})^{-1}
  \mathbf{E}^{\top}\mathbf{G}_{\mathrm{on}}.}{}
The only dense inverse is the small \eqn{r \times r}{} matrix
\eqn{(\mathbf{J}^{-1} + \mathbf{E}^{\top}\mathbf{G}_{\mathrm{on}}\mathbf{E})^{-1}}{}.

Define \eqn{\mathbf{N} = \mathbf{G}_{\mathrm{on}}^{1/2}\mathbf{E}}{} and a
thin QR decomposition \eqn{\mathbf{N} = \mathbf{Q}\mathbf{R}}{}. With
\eqn{\mathbf{P} = \mathbf{Q}\mathbf{Q}^{\top}}{} and
\eqn{\mathbf{S} = \mathbf{I}_r + \mathbf{R}\mathbf{J}\mathbf{R}^{\top}}{},
the supplement writes
\deqn{\mathbf{F}
  = (\mathbf{I}_P + \mathbf{N}\mathbf{J}\mathbf{N}^{\top})^{-1}
  = \mathbf{I}_P - \mathbf{P} + \mathbf{Q}\mathbf{S}^{-1}\mathbf{Q}^{\top},}{}
so that
\deqn{\mathbf{G}^{1/2}
  = \mathbf{G}_{\mathrm{on}}^{1/2}\mathbf{F}^{1/2}, \qquad
  \mathbf{G}^{-1/2}
  = \mathbf{F}^{-1/2}\mathbf{G}_{\mathrm{on}}^{-1/2},}{}
with
\deqn{\mathbf{F}^{1/2}
  = (\mathbf{I}_P - \mathbf{P}) +
  \mathbf{Q}\mathbf{S}^{-1/2}\mathbf{Q}^{\top}, \qquad
  \mathbf{F}^{-1/2}
  = (\mathbf{I}_P - \mathbf{P}) +
  \mathbf{Q}\mathbf{S}^{1/2}\mathbf{Q}^{\top}.}{}
This preserves the same transformed OLS projection used in the
independent case:
\deqn{\mathbf{y}^{*} = \mathbf{G}^{1/2}\mathbf{w}, \qquad
  \mathbf{X}^{*} = \mathbf{G}^{1/2}\mathbf{A}, \qquad
  \tilde{\boldsymbol{\beta}} = \mathbf{G}^{1/2}\mathbf{r}^{*},}{}
where \eqn{\mathbf{w} = \mathbf{X}^{\top}\mathbf{V}^{-1}\mathbf{y}}{}.
The full-space operations remain the QR on \eqn{\mathbf{X}^{*}}{} and a
rank-\eqn{r}{} correction through \eqn{\mathbf{Q}}{}.

\strong{Implementation map.} The code uses the same algebra but stores a
numerically convenient representation. The correspondence is:
\begin{itemize}

\item{} manuscript \eqn{\mathbf{E}}{} \eqn{\leftrightarrow}{} helper output
\code{E} (legacy alias \code{L});
\item{} manuscript diagonal sign matrix \eqn{\mathbf{J}}{}
\eqn{\leftrightarrow}{} helper output \code{J\_signs}
(legacy alias \code{S\_signs});
\item{} manuscript inverse
\eqn{(\mathbf{J}^{-1} + \mathbf{E}^{\top}\mathbf{G}_{\mathrm{on}}\mathbf{E})^{-1}}{}
\eqn{\leftrightarrow}{} helper output \code{inner\_inv}
(legacy alias \code{M});
\item{} manuscript projector basis \eqn{\mathbf{Q}}{}
\eqn{\leftrightarrow}{} internal basis \code{U\_Q}
(alias \code{Q\_basis});
\item{} manuscript projector \eqn{\mathbf{P} = \mathbf{Q}\mathbf{Q}^{\top}}{}
is not stored explicitly, but is represented through rank-\eqn{r}{}
multiplies involving \code{U\_Q};
\item{} manuscript factors \eqn{\mathbf{F}^{1/2}}{} and
\eqn{\mathbf{F}^{-1/2}}{} are stored internally as
\eqn{\mathbf{I} - \mathbf{U}_Q\mathbf{C}\mathbf{U}_Q^{\top}}{} and
\eqn{\mathbf{I} + \mathbf{U}_Q\mathbf{C}_{\mathrm{inv}}\mathbf{U}_Q^{\top}}{},
where \code{C} and \code{C\_inv} are diagonal.

\end{itemize}

These are just storage choices; the notation matches the supplement.

For the non-Gaussian Woodbury path
(\code{\LinkA{.get\_B\_gee\_glm\_woodbury}{.get.Rul.B.Rul.gee.Rul.glm.Rul.woodbury}}), the fixed correlation
perturbation \eqn{\mathbf{V}^{-1} - \mathbf{I}}{} and its product with
\eqn{\mathbf{X}}{} are precomputed once and held fixed
across Newton iterations. At each step, the weighted perturbation
\eqn{\mathbf{X}^{\top}\mathbf{W}(\mathbf{V}^{-1} -
\mathbf{I})\mathbf{X}}{} is formed using the
precomputed product, then split and re-factored via
\code{\LinkA{.woodbury\_redecompose\_weighted}{.woodbury.Rul.redecompose.Rul.weighted}}.

The Woodbury path is used when \eqn{r < P/3}{}; otherwise the code falls
back to the dense whitened solve. If \eqn{\mathbf{F}}{} is not positive
definite, the dense path is also used.
See \code{\LinkA{.get\_B\_gee\_woodbury}{.get.Rul.B.Rul.gee.Rul.woodbury}} (Gaussian) and
\code{\LinkA{.get\_B\_gee\_glm\_woodbury}{.get.Rul.B.Rul.gee.Rul.glm.Rul.woodbury}} (non-Gaussian).
\end{SubSection}


%
\begin{SubSection}{Step Control}
Different paths use different step-control strategies.

In the GEE paths (Paths 1a and 1b), the coefficient update is damped:
\deqn{\boldsymbol{\beta}^{(s+1)} = (1 - \alpha_s)\boldsymbol{\beta}^{(s)}
  + \alpha_s\boldsymbol{\beta}_{\mathrm{cand}}^{(s)},}{}
with \eqn{\alpha_s = 2^{-d_s}}{} and \eqn{d_s}{} increased whenever the
candidate gives non-finite or larger deviance. The loop terminates after
10 consecutive rejections or 100 total iterations, and exits early when
both coefficient change and deviance improvement fall below \code{tol}
after a burn-in period.

In Path 3 (non-Gaussian, no correlation), the outer projection loop has
no line search. The algorithm recomputes \eqn{\mathbf{G}}{} at the current
constrained estimate and applies a fresh projection. If the mean absolute
coefficient change begins increasing, the previous iterate is restored
and the loop stops. The partition-wise unconstrained estimates themselves
are obtained by damped Newton-Raphson inside
\code{\LinkA{unconstrained\_fit\_default}{unconstrained.Rul.fit.Rul.default}}, where
\code{\LinkA{damped\_newton\_r}{damped.Rul.newton.Rul.r}} computes a Newton direction once per
iteration and halves the step size until the penalized log-likelihood
improves.
\end{SubSection}

\end{Section}
%
\begin{Section}{Accommodating Correlation Structures}


%
\begin{SubSection}{Parametric Correlation Structures}
Suppose \eqn{\mathrm{Cov}(\mathbf{y}) = \sigma^2\,\mathbf{V}(\boldsymbol{\theta})}{}
for a known parametric family indexed by \eqn{\boldsymbol{\theta}}{} (e.g.,
AR(1) with \eqn{\theta = \rho}{}, Matern with
\eqn{\boldsymbol{\theta} = (\ell, \nu)}{}, exchangeable with
\eqn{\theta = \rho}{}). The penalized generalized least-squares problem
becomes
\deqn{\min_{\boldsymbol{\beta}}\;
  (\mathbf{y} - \mathbf{X}\boldsymbol{\beta})^{\top}\mathbf{V}^{-1}
  (\mathbf{y} - \mathbf{X}\boldsymbol{\beta})
  + \boldsymbol{\beta}^{\top}\boldsymbol{\Lambda}\boldsymbol{\beta}
  \quad \text{s.t.}\; \mathbf{A}^{\top}\boldsymbol{\beta} = \mathbf{0}.}{}
The correlation matrix \eqn{\mathbf{V}}{} is supplied through the
fitted-object components \code{Vhalf} and \code{VhalfInv}, either
directly or via user functions \code{Vhalf\_fxn} and \code{VhalfInv\_fxn}.
When both are non-\code{NULL}, \code{\LinkA{get\_B}{get.Rul.B}} dispatches to the GEE
paths (Path 1a or 1b). For built-in correlation structures, the required
square-root matrices are assembled numerically using
\code{\LinkA{matsqrt}{matsqrt}}, \code{\LinkA{matinvsqrt}{matinvsqrt}}, and \code{\LinkA{invert}{invert}}.
\end{SubSection}


%
\begin{SubSection}{Whitening and Permutation}
Because the data are stored in partition ordering (all observations from
partition 0, then partition 1, etc.) while \eqn{\mathbf{V}}{} is in the
original observation ordering, a permutation is applied internally:
\eqn{\mathbf{V}_{\mathrm{perm}}^{-1/2} =
\mathbf{V}^{-1/2}[\boldsymbol{\pi}, \boldsymbol{\pi}]}{}, where
\eqn{\boldsymbol{\pi} = \texttt{unlist(order\_list)}}{} maps original
indices to partition-ordered indices. The whitened design and response are
\eqn{\tilde{\mathbf{X}} =
\mathbf{V}_{\mathrm{perm}}^{-1/2}\mathbf{X}_{\mathrm{block}}}{} and
\eqn{\tilde{\mathbf{y}} =
\mathbf{V}_{\mathrm{perm}}^{-1/2}\mathbf{y}}{}.

The \eqn{\mathbf{X}}{} and \eqn{\mathbf{y}}{} inputs to
\code{\LinkA{lgspline.fit}{lgspline.fit}} are preserved in their unwhitened form.
Whitening is applied inside \code{\LinkA{get\_B}{get.Rul.B}} and
\code{\LinkA{blockfit\_solve}{blockfit.Rul.solve}} where the full \eqn{N \times P}{}
block-diagonal design is available, since applying
\eqn{\mathbf{V}^{-1/2}}{} to only the diagonal blocks of the partitioned
design would silently discard cross-partition contributions and corrupt
the Gram matrix.
\end{SubSection}


%
\begin{SubSection}{Loss of Block-Diagonal Structure}
Unlike the independent-errors case,
\eqn{\mathbf{X}^{\top}\mathbf{V}^{-1}\mathbf{X}}{} is not block-diagonal
because \eqn{\mathbf{V}^{-1}}{} introduces cross-partition coupling. The
unconstrained estimator
\deqn{\hat{\boldsymbol{\beta}} =
  (\mathbf{X}^{\top}\mathbf{V}^{-1}\mathbf{X} +
  \boldsymbol{\Lambda})^{-1}
  \mathbf{X}^{\top}\mathbf{V}^{-1}\mathbf{y}}{}
requires a full \eqn{P \times P}{} solve, and the constraint projection
proceeds with
\eqn{\mathbf{G} = (\mathbf{X}^{\top}\mathbf{V}^{-1}\mathbf{X} +
\boldsymbol{\Lambda})^{-1}}{}.

For structured correlation matrices (AR(1), exchangeable, banded), the
perturbation \eqn{\mathbf{V}^{-1} - \mathbf{I}}{} is sparse and the
cross-partition coupling has low effective rank. In these cases the
Woodbury-accelerated paths
(\code{\LinkA{.get\_B\_gee\_woodbury}{.get.Rul.B.Rul.gee.Rul.woodbury}},
\code{\LinkA{.get\_B\_gee\_glm\_woodbury}{.get.Rul.B.Rul.gee.Rul.glm.Rul.woodbury}}) recover the partition-wise
computational structure by decomposing the coupling into a
block-diagonal correction plus a low-rank remainder, as described in the
GLM Extension section. For dense or high-rank
\eqn{\mathbf{V}^{-1}}{}, the code falls back to the full whitened system
(\code{\LinkA{.get\_B\_gee\_gaussian}{.get.Rul.B.Rul.gee.Rul.gaussian}}, \code{\LinkA{.get\_B\_gee\_glm}{.get.Rul.B.Rul.gee.Rul.glm}}).
\end{SubSection}


%
\begin{SubSection}{GEE Deviance Monitoring}
For non-Gaussian models with correlation, the deviance used for
convergence monitoring is computed in the whitened space by
\code{.bf\_gee\_deviance}. When the family supplies
\code{custom\_dev.resids}, the raw deviance residuals
\eqn{r_i = \mathrm{sign}(d_i)\sqrt{|d_i|}}{} are divided by
\eqn{\sqrt{w_i}}{} and pre=dmultiplied by
\eqn{\mathbf{V}_{\mathrm{perm}}^{-1/2}}{} before squaring and averaging:
\deqn{D_{\mathrm{GEE}} = \frac{1}{N}\left\|
  \mathbf{V}_{\mathrm{perm}}^{-1/2}\,
  \mathbf{w}^{-1/2}\mathbf{r}\right\|^{2},}{}
where \eqn{\mathbf{w}}{} is the vector of working weights at the current
iterate, clamped below at \eqn{\sqrt{\varepsilon_{\mathrm{mach}}}}{}. When
only \code{dev.resids} is available, the function falls back to the
standard mean deviance; otherwise it uses mean squared error.
\end{SubSection}


%
\begin{SubSection}{REML Estimation of Correlation Parameters}
Correlation parameters \eqn{\boldsymbol{\theta}}{} are estimated by minimizing
a negative restricted log-likelihood (REML) objective. The criterion
implemented in \pkg{lgspline} is a central-limit-theorem-based working
approximation to a Laplace-style marginal likelihood criterion,
applied here solely to correlation structure estimation rather than
penalty parameter selection.
Let \eqn{\mathbf{D} = \mathrm{diag}(d_i)}{} be the observation weight matrix,
\eqn{\tilde{\mathbf{W}} = \mathrm{diag}(\tilde{w}_i)}{} the GLM working weight
matrix at the current fitted values, \eqn{\mathbf{V}}{} the correlation matrix
parameterized by \eqn{\boldsymbol{\rho}}{} (a vector on the unconstrained
real line), and \eqn{\tilde{\sigma}^{2}}{} the dispersion profiled at its
restricted maximum likelihood estimate. The negative REML objective
implemented in \pkg{lgspline}, scaled by \eqn{1/N}{}, is
\deqn{-\ell_{R}(\boldsymbol{\rho}) = \frac{1}{N}\!\left[
  -\log|\mathbf{V}^{-1/2}|
  + \frac{N}{2}\log\tilde{\sigma}^{2}
  + \frac{1}{2\tilde{\sigma}^{2}}
    (\mathbf{y} - \boldsymbol{\mu})^{\top}
    \mathbf{D}\tilde{\mathbf{W}}^{-1}\mathbf{V}^{-1}
    (\mathbf{y} - \boldsymbol{\mu})
  + \frac{1}{2}\log|(\tilde{\sigma}^{2}\mathbf{U}\mathbf{G})^{-1}|^{+}
\right],}{}
where \eqn{|\cdot|^{+}}{} denotes the generalized determinant (product of
nonzero eigenvalues), and
\eqn{\boldsymbol{\mu} = g^{-1}(\mathbf{X}\tilde{\boldsymbol{\beta}})}{} are
the fitted values on the response scale.

Gradients with respect to correlation parameters are available in
closed form for all built-in structures except Matern, which uses
finite-difference approximation due to the complexity of differentiating
the modified Bessel function \eqn{K_{\nu}}{} with respect to \eqn{\nu}{}. See
\code{\LinkA{reml\_grad\_from\_dV}{reml.Rul.grad.Rul.from.Rul.dV}} for the full gradient derivation and
notation. Custom analytic gradients can be supplied through \code{REML\_grad},
and a fully custom criterion can replace REML through
\code{custom\_VhalfInv\_loss}. The Toeplitz example in \code{\LinkA{lgspline}{lgspline}}
demonstrates how to supply custom correlation structures with user-defined
gradient functions. Optimization over these working correlation parameters is
then carried out by the same quasi-Newton engine used elsewhere in the
package, with a finite-difference fallback when needed. In the user-facing interface, this
machinery is activated through \code{correlation\_structure},
\code{correlation\_id}, \code{spacetime}, \code{VhalfInv}, \code{Vhalf},
\code{VhalfInv\_fxn}, \code{Vhalf\_fxn}, \code{VhalfInv\_par\_init},
\code{REML\_grad}, \code{custom\_VhalfInv\_loss}, and
\code{VhalfInv\_logdet}.

The gradient of the negative REML has three terms per parameter:
\begin{enumerate}

\item{} \eqn{\frac{1}{2}\mathrm{tr}(\mathbf{V}^{-1}\partial\mathbf{V}/\partial\theta_j)}{}:
the log-determinant contribution.
\item{} \eqn{-\frac{1}{2\tilde{\sigma}^{2}}\mathbf{r}^{\top}
    (\partial\mathbf{V}/\partial\theta_j)\mathbf{r}}{}:
the residual quadratic form contribution, where
\eqn{\mathbf{r} = \mathrm{diag}(\sqrt{d_i}/\sqrt{\tilde{w}_i})
    \mathbf{V}^{-1/2}(\mathbf{y} - \boldsymbol{\mu})}{}.
\item{} \eqn{-\frac{1}{2}\mathrm{tr}\!\left(\mathbf{M}^{+}\mathbf{X}_{*}^{\top}
    \mathbf{V}^{-1}(\partial\mathbf{V}/\partial\theta_j)\mathbf{V}^{-1}
    \tilde{\mathbf{W}}\mathbf{D}\mathbf{X}_{*}\right)}{}:
the REML correction, where
\eqn{\mathbf{X}_{*} = \mathbf{X}\mathbf{U}}{} is the constrained design and
\eqn{\mathbf{M} = \mathbf{X}_{*}^{\top}\mathbf{V}^{-1}
    \tilde{\mathbf{W}}\mathbf{D}\mathbf{X}_{*}
    + \mathbf{U}^{\top}\boldsymbol{\Lambda}\mathbf{U}}{} is the projected
penalized information.

\end{enumerate}


For each supported correlation family, the derivatives
\eqn{\partial\mathbf{V}/\partial\theta_j}{} are available in closed form,
enabling analytic gradient computation for use with the quasi-Newton
optimizer.
\end{SubSection}


%
\begin{SubSection}{Connection to Standard Mixed Model REML}
The quadratic penalty \eqn{\boldsymbol{\Lambda}}{} acts as the inverse prior
covariance of a Gaussian random effect on the spline coefficients, with
the smoothing parameter satisfying \eqn{\lambda = \tau^{2}/\sigma^{2}}{};
this is the same mixed model representation used by \pkg{mgcv}, and
Monte Carlo draws under the resulting Laplace-style approximation are
available via \code{\LinkA{generate\_posterior}{generate.Rul.posterior}}. For Gaussian responses
with identity link and no penalization, the implemented criterion
coincides exactly with classical REML. For non-Gaussian responses, the
criterion substitutes the Fisher information for the full penalized
log-likelihood Hessian, exploiting the CLT approximation that
\eqn{\tilde{\mathbf{W}}^{-1/2}(\mathbf{y} - \boldsymbol{\mu})}{} is
approximately Gaussian; this yields a method-of-moments style estimator
for the correlation and variance parameters that depends only on
mean-variance relationships through \eqn{\mathbf{W}}{}, and therefore
generalizes naturally to quasi-likelihood and other settings where a
fully specified log-likelihood is unavailable. The REML correction term
\eqn{\log|(\tilde{\sigma}^{2}\mathbf{U}\mathbf{G})^{-1}|^{+}}{} uses a
generalized log-determinant so that only nonzero eigenvalues contribute
when rank deficiency arises from smoothness constraints or
identifiability conditions in \eqn{\mathbf{A}}{}. When
\eqn{\boldsymbol{\Lambda}}{} is full rank, the criterion coincides with
the exact marginal likelihood from integrating out
\eqn{\boldsymbol{\beta}}{} under its Gaussian prior; when
\eqn{\boldsymbol{\Lambda}}{} is rank-deficient, unpenalized coefficients
are projected out in the REML sense while penalized coefficients are
integrated through their prior. During penalty tuning, the block-diagonal
approximation is retained for GCV criteria and gradients; since GCV is
rotation-invariant the practical effect on automatic selection is expected
to be negligible, though this is not confirmed and the tuned penalties can
always be overridden.
\end{SubSection}


%
\begin{SubSection}{Built-In Correlation Structures}
The package provides several built-in correlation structures for modeling
spatial and temporal dependence. These are specified via
\code{correlation\_structure} with group membership in \code{correlation\_id}
and spatial or temporal coordinates in \code{spacetime} (an \eqn{N}{}-row
matrix). Exchangeable correlation does not require \code{spacetime}.

All positive scale parameters are estimated on the log scale,
with back-transform \eqn{\exp(\cdot)}{}. Parameters constrained to
\eqn{(0, 1)}{} use a double-exponential back-transform of the form
\eqn{\exp(-\exp(\eta))}{}, so optimization still occurs on the
unconstrained real line while the correlation remains bounded.

\begin{description}

\item[Exchangeable] 
Aliases: \code{'exchangeable'}, \code{'cs'}, \code{'CS'},
\code{'compoundsymmetric'}, \code{'compound-symmetric'}.

A constant correlation \eqn{\nu}{} between any two observations within
the same cluster. Parameterization: \eqn{\nu = \exp(-\exp(\rho))}{},
so \eqn{\nu \in (0, 1)}{}. Only positive within-cluster correlation is
supported under this parameterization.

\item[Spatial Exponential] 
Aliases: \code{'spatial-exponential'}, \code{'spatialexponential'},
\code{'exp'}, \code{'exponential'}.

Correlation decays exponentially with distance:
\eqn{\exp(-\omega d)}{} where \eqn{d}{} is Euclidean distance and
\eqn{\omega > 0}{}. Parameterization: \eqn{\omega = \exp(\rho)}{}.
Mathematically equivalent to the power correlation \eqn{\theta^{d}}{}
with \eqn{\theta = e^{-\omega}}{}, but with better numerical properties
during optimization.

\item[AR(1)] 
Aliases: \code{'ar1'}, \code{'ar(1)'}, \code{'AR(1)'}, \code{'AR1'}.

Correlation depends on rank difference between observations:
\eqn{\nu^{r}}{} where \eqn{r}{} is the rank difference within cluster.
Parameterization: \eqn{\nu = \exp(-\exp(\rho))}{},
so \eqn{\nu \in (0, 1)}{}. Only positive autocorrelation is supported.

\item[Gaussian / Squared Exponential] 
Aliases: \code{'gaussian'}, \code{'rbf'}, \code{'squared-exponential'}.

Smooth decay with squared distance:
\eqn{\exp(-d^{2}/(2\ell^{2}))}{} where \eqn{\ell}{} is the length scale.
Parameterization: \eqn{\ell = \exp(\rho)}{}.

\item[Spherical] 
Aliases: \code{'spherical'}, \code{'Spherical'}, \code{'cubic'},
\code{'sphere'}.

Polynomial decay with a hard cutoff at range \eqn{r}{}:
\eqn{1 - 1.5(d/r) + 0.5(d/r)^{3}}{} for \eqn{d \le r}{}, and \eqn{0}{}
otherwise. Parameterization: \eqn{r = \exp(\rho)}{}.

\item[Matern] 
Aliases: \code{'matern'}, \code{'Matern'}.

Flexible correlation with adjustable smoothness:
\eqn{(2^{1-\nu}/\Gamma(\nu))(\sqrt{2\nu}\,d/\ell)^{\nu}K_{\nu}(\sqrt{2\nu}\,d/\ell)}{}.
Two parameters: length scale \eqn{\ell = \exp(\rho_1)}{} and smoothness
\eqn{\nu = \exp(\rho_2)}{}. No analytical gradient is available for
\eqn{\nu}{} due to the difficulty of differentiating the modified Bessel
function \eqn{K_{\nu}}{} with respect to its order, so finite differences
are used; this makes Matern slower and potentially less stable than
other structures.

\item[Gamma-Cosine] 
Aliases: \code{'gamma-cosine'}, \code{'gammacosine'},
\code{'GammaCosine'}.

Oscillatory dependence:
\eqn{(d^{\alpha-1}e^{-\gamma d})/(\Gamma(\alpha)/\gamma^{\alpha})\cdot\cos(\omega d)}{}.
Three parameters: shape \eqn{\alpha = \exp(\rho_1)}{}, rate
\eqn{\gamma = \exp(\rho_2)}{}, frequency \eqn{\omega = \exp(\rho_3)}{}.
Reduces to exponential when \eqn{\alpha = 1}{} and \eqn{\omega \approx 0}{}.

\item[Gaussian-Cosine] 
Aliases: \code{'gaussian-cosine'}, \code{'gaussiancosine'},
\code{'GaussianCosine'}.

Smooth oscillatory correlation:
\eqn{\exp(-d^{2}/(2\ell^{2}))\cdot\cos(\omega d)}{}.
Two parameters: length scale \eqn{\ell = \exp(\rho_1)}{} and frequency
\eqn{\omega = \exp(\rho_2)}{}. Reduces to Gaussian when
\eqn{\omega \approx 0}{}.


\end{description}

\end{SubSection}


%
\begin{SubSection}{Interpreting Estimated Correlation Parameters}
Correlation parameters are estimated on transformed scales; they must be
back-transformed for interpretation. When \code{\LinkA{confint.lgspline}{confint.lgspline}} is called and the
inverse Hessian from BFGS is available, confidence intervals are returned
on the untransformed (working) scale and should be back-transformed as
described in the examples for \code{\LinkA{lgspline}{lgspline}}.
\end{SubSection}


%
\begin{SubSection}{Custom Correlation Structures}
Custom correlation structures can be specified through:
\begin{itemize}

\item{} \code{VhalfInv\_fxn}: Creates \eqn{\mathbf{V}^{-1/2}}{}.
\item{} \code{Vhalf\_fxn}: Creates \eqn{\mathbf{V}^{1/2}}{}. When omitted,
the code computes it by explicit inversion of \code{VhalfInv}.
\item{} \code{REML\_grad}: Provides the analytical gradient of the REML
objective.
\item{} \code{VhalfInv\_logdet}: Efficient log-determinant computation.
\item{} \code{custom\_VhalfInv\_loss}: Replaces the REML objective entirely.

\end{itemize}

These functions enter \code{\LinkA{lgspline}{lgspline}} through
\code{correlation\_structure}, \code{VhalfInv\_fxn}, \code{Vhalf\_fxn},
\code{REML\_grad}, and \code{custom\_VhalfInv\_loss}, and the fitted object
retains the resulting correlation machinery in components such as
\code{VhalfInv\_fxn}, \code{Vhalf\_fxn}, and
\code{VhalfInv\_params\_estimates}. When \code{VhalfInv} is supplied but
\code{Vhalf} is not, \code{Vhalf} is computed unconditionally as the inverse
of \code{VhalfInv} for all family/link combinations, since both
\code{\LinkA{get\_B}{get.Rul.B}} and \code{\LinkA{blockfit\_solve}{blockfit.Rul.solve}} require it for GEE
estimation.
\end{SubSection}

\end{Section}
%
\begin{Section}{Variance and Dispersion Estimation}


Once the constrained estimate has been obtained, the next questions are how
much flexibility the fitted model effectively used and how uncertainty should
be propagated through the same constrained geometry. The quantities in this
section are therefore all built on the projected information matrices from
the previous sections.

%
\begin{SubSection}{Effective Degrees of Freedom and Dispersion}
The effective degrees of freedom is the trace of the hat matrix. In the
Gaussian identity case with observation weights and correlation, the fitted
linear operator is built from the dense GLS analogue
\eqn{\mathbf{G}_{\mathrm{correct}}}{} and can be written schematically as
\deqn{\mathbf{H} =
  \mathbf{V}^{-1/2}(\mathbf{W}\mathbf{D})^{1/2}
  \mathbf{X}\mathbf{U}\mathbf{G}_{\mathrm{correct}}\mathbf{X}^{\top}
  (\mathbf{W}\mathbf{D})^{1/2}\mathbf{V}^{-1/2},}{}
where for Gaussian identity \eqn{\mathbf{W} = \mathbf{I}}{}. In the
no-correlation Gaussian case this reduces to the familiar
\eqn{\mathbf{H} = \mathbf{X}\mathbf{U}\mathbf{G}\mathbf{X}^{\top}\mathbf{D}}{}.

For Gaussian identity fits, the dispersion estimate is computed as a
weighted mean squared residual, optionally scaled by
\eqn{N/(N - \mathrm{tr}(\mathbf{H}))}{} when
\code{unbias\_dispersion = TRUE}:
\deqn{\tilde{\sigma}^{2} =
  \frac{1}{N - \mathrm{tr}(\mathbf{H})}\|\mathbf{y} - \mathbf{\tilde{y}}_i \|^{2}.}{}

More generally with weights, a correlation structure and non-linear link function:
\deqn{\tilde{\sigma}^{2} =
  \frac{1}{N - \mathrm{tr}(\mathbf{H})}\| \mathbf{V}^{-1/2}\mathbf{W}^{-1/2}\mathbf{D}^{1/2}(\mathbf{y} - \tilde{\mathbf{y}})\|^{2}.}{}

This estimated dispersion is returned as \code{sigmasq\_tilde}, and the
corresponding effective degrees of freedom trace is returned as
\code{trace\_XUGX}. For non-Gaussian families, the fitting code delegates
dispersion estimation to \code{dispersion\_function}; thus the package does
not assume a single closed-form Pearson-style formula outside the Gaussian
identity setting. The hat-matrix trace itself is assembled by
\code{\LinkA{compute\_trace\_H}{compute.Rul.trace.Rul.H}} in the dense correlation-aware case and by the
same blockwise products summarized by \code{trace\_XUGX} in the simpler
no-correlation paths.
A concrete built-in non-Gaussian example is the Weibull AFT path, which pairs
\code{\LinkA{weibull\_family}{weibull.Rul.family}} with \code{\LinkA{weibull\_dispersion\_function}{weibull.Rul.dispersion.Rul.function}},
\code{\LinkA{weibull\_glm\_weight\_function}{weibull.Rul.glm.Rul.weight.Rul.function}}, and
\code{\LinkA{weibull\_schur\_correction}{weibull.Rul.schur.Rul.correction}}.
Users who want these quantities available for downstream inference should
keep \code{estimate\_dispersion = TRUE} and \code{return\_varcovmat = TRUE}
(the defaults), since \code{\LinkA{wald\_univariate}{wald.Rul.univariate}},
\code{\LinkA{confint.lgspline}{confint.lgspline}}, and the prediction-standard-error path in
\code{\LinkA{predict.lgspline}{predict.lgspline}} all rely on the post-fit dispersion and
covariance components documented here.
\end{SubSection}


%
\begin{SubSection}{Variance-Covariance Matrix}
The variance-covariance matrix of \eqn{\tilde{\boldsymbol{\beta}}}{} is
estimated as:
\deqn{\mathrm{Var}(\tilde{\boldsymbol{\beta}}) = \tilde{\sigma}^{2}(\mathbf{U}\mathbf{G}^{1/2})(\mathbf{U}\mathbf{G}^{1/2})^{\top}}{}
using the outer-product form for numerical stability. The result is returned
as \code{varcovmat} when \code{return\_varcovmat = TRUE}. The algebraically
equivalent expression \eqn{\tilde{\sigma}^{2}\mathbf{U}\mathbf{G}}{} is not
used because \eqn{\mathbf{G}}{} is only positive semi-definite when the
penalty matrix \eqn{\boldsymbol{\Lambda}}{} has zero eigenvalues (e.g., the
intercept and linear terms under the smoothing spline penalty when
\code{flat\_ridge\_penalty} = 0), which can introduce negative diagonal
entries in finite precision arithmetic. The outer-product form also
guarantees symmetry.

This is the Bayesian posterior covariance, treating the penalty as a
Gaussian prior on the coefficients. When \code{exact\_varcovmat = TRUE},
a frequentist correction is additionally computed:
\deqn{\mathrm{Var}_{\mathrm{exact}}(\tilde{\boldsymbol{\beta}}) =
  \tilde{\sigma^2}\mathbf{U}\mathbf{G}^{1/2}(\mathbf{X}^{\top}\mathbf{W}\mathbf{D}\mathbf{V}^{-1}\mathbf{X})\mathbf{G}^{1/2}\mathbf{U}^{\top} =
  \tilde{\sigma}^{2}\mathbf{U}\mathbf{G}\mathbf{U}^{\top}
  - \tilde{\sigma}^{2}\mathbf{U}\mathbf{G}\boldsymbol{\Lambda}\mathbf{G}\mathbf{U}^{\top}.}{}
The first term is the Bayesian posterior covariance; the second is a frequentist
correction such that for Gaussian identity link (with or without correlation),
this is the exact variance-covariance matrix of the constrained estimator.

When a correlation structure is present (\code{VhalfInv} non-\code{NULL}),
the block-diagonal \eqn{\mathbf{G}}{} is replaced by the full weighted GLS
analogue
\deqn{\mathbf{G}_{\mathrm{correct}} =
  \left(\mathbf{X}^{\top}\mathbf{W}\mathbf{D}\mathbf{V}^{-1}\mathbf{X}
  + \boldsymbol{\Lambda}\right)^{-1},}{}
where \eqn{\mathbf{W} = \mathbf{I}}{} in the Gaussian identity case.
This dense matrix is what enters the correlation-aware \eqn{\mathbf{U}}{},
\eqn{\mathrm{Var}(\tilde{\boldsymbol{\beta}})}{}, and
\eqn{\mathrm{Var}_{\mathrm{exact}}(\tilde{\boldsymbol{\beta}})}{}
computations.

In user-facing terms, \code{return\_varcovmat} controls whether this matrix is
stored at all, while \code{exact\_varcovmat} switches between the default
posterior/Laplace approximation and the exact frequentist correction in the
Gaussian-identity setting. The stored covariance is what powers
\code{\LinkA{wald\_univariate}{wald.Rul.univariate}}, \code{\LinkA{confint.lgspline}{confint.lgspline}}, and
\code{se.fit = TRUE} in \code{\LinkA{predict.lgspline}{predict.lgspline}}; the
\code{critical\_value} argument supplied at fit time is carried forward as the
default cutoff for those interval-producing helpers.
\end{SubSection}


%
\begin{SubSection}{Recomputation of G at Convergence}
At the final iterate, \eqn{\mathbf{G}}{} is recomputed to reflect the
converged working weights and Schur corrections. The implementation
computes the weighted design
\eqn{\mathbf{X}_{w}^{(k)} = \mathbf{X}_k \cdot \mathrm{diag}(\sqrt{\mathbf{w}_k})}{},
forms the weighted Gram matrix
\eqn{\mathbf{X}_{w}^{(k)\top}\mathbf{X}_{w}^{(k)}}{}, adds the Schur
correction, and performs eigendecomposition via
\code{\LinkA{compute\_G\_eigen}{compute.Rul.G.Rul.eigen}} to obtain \eqn{\mathbf{G}_k}{},
\eqn{\mathbf{G}_k^{1/2}}{}, and \eqn{\mathbf{G}_k^{-1/2}}{}. The relationship
\eqn{\mathbf{G}_k = \mathbf{G}_k^{1/2}(\mathbf{G}_k^{1/2})^{\top}}{} is
enforced exactly by construction, and the fitted object can retain these as
\code{G} and \code{Ghalf} when \code{return\_G = TRUE} and
\code{return\_Ghalf = TRUE}. The square-root factors are numerically
stabilized through the helper routines \code{\LinkA{matsqrt}{matsqrt}} and
\code{\LinkA{matinvsqrt}{matinvsqrt}}, which are also used elsewhere in the package when
dense analogues of \eqn{\mathbf{G}^{1/2}}{} or \eqn{\mathbf{G}^{-1/2}}{} are
required.
\end{SubSection}

\end{Section}
%
\begin{Section}{Bayesian Interpretation}


The penalty has a natural Gaussian-prior interpretation, so once the
constrained estimator and its covariance are available, Bayesian-style
posterior simulation follows almost immediately. This section records the
interpretation that is already implicit in the fitted object and in the
package's posterior simulation helpers.

A Bayesian interpretation follows from viewing the penalty as a Gaussian
prior on the coefficients. Conditional on the fitted smoothing parameters,
the code samples on the coefficient scale from
\deqn{\boldsymbol{\beta}^{(m)} =
  \tilde{\boldsymbol{\beta}} +
  \sqrt{\tilde{\sigma}^{2}},\mathbf{U}\mathbf{G}^{1/2}\mathbf{z}^{(m)},
  \qquad \mathbf{z}^{(m)} \sim \mathcal{N}(\mathbf{0}, \mathbf{I}).}{}
The coefficients are then back-transformed to the original response and
predictor scales.

When inequality constraints are absent, these draws are i.i.d. Gaussian
posterior draws around the fitted mode.
The underlying coefficient-draw closure also contains an elliptical slice
sampling route for active inequality constraints, using the same covariance
factor to keep retained draws in the feasible region.

At the implementation level, standard Gaussian posterior draws may place
positive mass on the infeasible region, so the constrained-draw closure
instead targets the truncated posterior
\deqn{\pi(\boldsymbol{\beta} \mid \mathbf{y}) \propto
  \exp\!\left(-\frac{1}{2}(\boldsymbol{\beta} - \tilde{\boldsymbol{\beta}})^{\top}
  \mathbf{G}^{-1}(\boldsymbol{\beta} - \tilde{\boldsymbol{\beta}})\right)
  \mathbf{1}(\mathbf{C}^{\top}\boldsymbol{\beta} \succeq \mathbf{c}),}{}
yielding credible intervals that respect the constraint boundaries.
The public \code{\LinkA{generate\_posterior}{generate.Rul.posterior}} wrapper forwards
\code{enforce\_qp\_constraints} to the stored constrained-draw closure, so
constrained draws can be requested directly from the user-facing interface.
When a working correlation structure is present, the companion helper
\code{\LinkA{generate\_posterior\_correlation}{generate.Rul.posterior.Rul.correlation}} extends this idea by propagating
uncertainty in the fitted correlation parameters through the same
\code{VhalfInv\_fxn}/\code{Vhalf\_fxn} machinery described in the correlation
section, rather than conditioning only on fixed covariance parameters. Correlation
parameters are drawn from a multivariate normal distribution centered about their estimates
with the inverse approxiamte BFGS Hessian of the REML optimization problem
\code{VhalfInv\_params\_vcov} used by default (or a custom alternative as
supplied to the argument \code{correlation\_param\_vcov\_sc}).
\end{Section}
%
\begin{Section}{Inequality Constraints via Sequential Quadratic Programming}


%
\begin{SubSection}{Overview}
Inequality constraints of the form
\eqn{\mathbf{C}^{\top}\boldsymbol{\beta} \succeq \mathbf{c}}{} handle
shape restrictions such as monotonicity, convexity, or boundedness. In
the monomial basis, these are linear in \eqn{\boldsymbol{\beta}}{}.
Monotonicity at a grid of points requires the first derivative polynomial
to be non-negative there; convexity requires the second derivative to be
non-negative; range constraints bound the fitted values directly. All of
these translate to linear inequality constraints on the coefficient
vector.

The inequality pieces are assembled by \code{\LinkA{process\_qp}{process.Rul.qp}}, which
returns the \code{qp\_Amat}, \code{qp\_bvec}, and \code{qp\_meq} objects
passed to \code{\LinkA{solve.QP}{solve.QP}}, together with a
\code{quadprog} flag indicating whether any inequality constraints are
active. For derivative-sign constraints, \code{\LinkA{process\_qp}{process.Rul.qp}} calls
\code{.build\_deriv\_qp}, which uses
\code{\LinkA{make\_derivative\_matrix}{make.Rul.derivative.Rul.matrix}} on the expansion-standardized
design and maps the derivative rows into the full \eqn{P}{}-dimensional
coefficient space partition by partition.
\end{SubSection}


%
\begin{SubSection}{Partition-Wise Active-Set Method}
The sparsity pattern of \eqn{\mathbf{C}}{} is inspected automatically by
\code{\LinkA{.detect\_qp\_global}{.detect.Rul.qp.Rul.global}} (equivalently
\code{.solver\_detect\_qp\_global}). When every column of
\eqn{\mathbf{C}}{} has nonzero entries in only a single partition block,
the constraint system is block-separable and an active-set method can
replace the dense QP.

At each iteration, the active set \eqn{\mathcal{A}}{} (constraints
satisfied at equality) is appended to \eqn{\mathbf{A}}{} as additional
equality constraints:
\deqn{\mathbf{A}_{\mathrm{aug}} = [\mathbf{A} \mid
  \mathbf{C}_{\mathcal{A}}].}{}
The constrained estimate is obtained by the same OLS projection with
\eqn{\mathbf{A}_{\mathrm{aug}}}{} in place of \eqn{\mathbf{A}}{}, and
since \eqn{\mathbf{A}_{\mathrm{aug}}}{} retains block-diagonal
compatibility, all operations remain partition-wise.

The active set is updated by checking primal feasibility (adding the
most-violated inactive constraint) and dual feasibility (dropping the
active constraint with the most negative Lagrange multiplier) until the
KKT conditions are satisfied. Lagrange multipliers for active
inequalities are recovered from the OLS fit used in the Lagrangian
projection: the fitted coefficients on
\eqn{\mathbf{X}^* = \mathbf{G}^{1/2}\mathbf{A}_{\mathrm{aug}}}{} give
the multipliers up to sign. Implemented in
\code{\LinkA{.active\_set\_refine}{.active.Rul.set.Rul.refine}} and
\code{\LinkA{.check\_kkt\_partitionwise}{.check.Rul.kkt.Rul.partitionwise}}.

When any column of \eqn{\mathbf{C}}{} spans multiple partition blocks
(for example, cross-knot monotonicity constraints), the block-diagonal
structure is broken and the dense SQP approach is required. The
selection is made automatically. If the active-set method does not
converge within its iteration limit (default 50), the code falls back
to dense SQP as well.
\end{SubSection}


%
\begin{SubSection}{Dense SQP Iteration}
The dense SQP approach, implemented in \code{\LinkA{.qp\_refine}{.qp.Rul.refine}},
solves a sequence of quadratic subproblems approximating the penalized
log-likelihood. At each iteration \eqn{s}{}:
\begin{enumerate}

\item{} Compute the information matrix
\eqn{\mathbf{M}^{(s)} =
    \mathbf{X}^{\top}\mathbf{W}^{(s)}\mathbf{X} +
    \boldsymbol{\Lambda}_{\mathrm{block}} + \mathbf{S}^{(s)}}{},
where \eqn{\mathbf{S}^{(s)}}{} is the Schur complement correction.
\item{} Compute the score vector via \code{qp\_score\_function}.
\item{} Solve the QP with \code{\LinkA{solve.QP}{solve.QP}}, where the
combined constraint matrix is
\eqn{[\mathbf{A} \mid \mathbf{C}]}{} with the first \eqn{R}{} columns
treated as equalities.
\item{} Apply a damped update
\eqn{\boldsymbol{\beta}^{(s+1)} =
    (1 - \alpha)\boldsymbol{\beta}^{(s)} +
    \alpha\boldsymbol{\beta}_{\mathrm{QP}}}{}, with
\eqn{\alpha = 2^{-d}}{} and \eqn{d}{} incremented upon deviance
increase.

\end{enumerate}

A rescaling factor
\eqn{\mathrm{sc} = \sqrt{\mathrm{mean}(|\mathbf{M}^{(s)}|)}}{} is applied
to the Hessian and linear term before calling the QP solver. The
equality-constrained estimate from the Lagrangian projection serves as
a warm start. The \code{qp\_score\_function} defaults to the canonical GLM
score
\eqn{\mathbf{X}^{\top}(\mathbf{y} - \boldsymbol{\mu})}{}; for custom
models a different score can be supplied.
\end{SubSection}


%
\begin{SubSection}{Active Set and Lagrange Multipliers}
The active set at the solution identifies binding inequality constraints.
The Lagrange multipliers quantify the cost of each: a multiplier of zero
means the constraint is not binding. The implementation stores active
constraint indices, the corresponding submatrix of the constraint matrix,
and the multiplier vector in the \code{qp\_info} list returned alongside
coefficient estimates, with components \code{lagrangian}, \code{iact},
and \code{Amat\_active}. When
\code{return\_lagrange\_multipliers = TRUE}, the fitted object also stores
the final multiplier vector directly as \code{lagrange\_multipliers}.

The original assembled inequality data are retained in the fitted
object's \code{quadprog\_list} component (containing \code{qp\_Amat},
\code{qp\_bvec}, and \code{qp\_meq} from \code{\LinkA{process\_qp}{process.Rul.qp}}) so
the final active set can be interpreted relative to the full
specification. The final \eqn{\mathbf{U}}{} used in constructing the
posterior variance-covariance matrix is built from both the equality
constraints and the active inequality constraints at the solution.
\end{SubSection}


%
\begin{SubSection}{Built-In Constraints}
Built-in inequality constraints include:
\begin{itemize}

\item{} Monotonicity: \code{qp\_monotonic\_increase},
\code{qp\_monotonic\_decrease}. Enforced by requiring consecutive
fitted values to be non-decreasing (or non-increasing):
\eqn{(\mathbf{x}_i - \mathbf{x}_{i-1})^{\top}\boldsymbol{\beta}
    \geq 0}{}. These are constructed by \code{\LinkA{process\_qp}{process.Rul.qp}} from
the partition-stacked block design reordered to observation order.
\item{} Derivative sign: \code{qp\_positive\_derivative},
\code{qp\_negative\_derivative}. Enforced through the
first-derivative design matrix from
\code{\LinkA{make\_derivative\_matrix}{make.Rul.derivative.Rul.matrix}}. May be \code{TRUE}/\code{FALSE}
or a character/integer vector selecting specific predictors.
\item{} Second-derivative sign: \code{qp\_positive\_2ndderivative},
\code{qp\_negative\_2ndderivative}. Same construction using the
second-derivative design matrix.
\item{} Response range: \code{qp\_range\_lower}, \code{qp\_range\_upper}.
Bounds on the linear predictor; for non-identity links, the bounds
are transformed to the link scale inside \code{\LinkA{process\_qp}{process.Rul.qp}}.
\item{} Custom constraints via \code{qp\_Amat\_fxn}, \code{qp\_bvec\_fxn},
and \code{qp\_meq\_fxn}, which receive the design matrix structure
and return the constraint matrix, bound vector, and number of
equalities. These are commonly paired with a custom
\code{qp\_score\_function} when the quadratic approximation uses a
non-default likelihood. The low-level objects \code{qp\_Amat},
\code{qp\_bvec}, and \code{qp\_meq} remain documented in
\code{\LinkA{lgspline}{lgspline}} but in the current implementation serve as
activation markers rather than being merged into the constraint set
assembled by \code{\LinkA{process\_qp}{process.Rul.qp}}.

\end{itemize}

All constraints can be thinned to a user-specified subset of rows via
\code{qp\_observations}, which \code{\LinkA{process\_qp}{process.Rul.qp}} applies before
assembly.
\end{SubSection}

\end{Section}
%
\begin{Section}{Blockfit Backfitting for Linear Non-Interactive Effects}


%
\begin{SubSection}{Motivation}
When a model contains both spline terms (receiving \eqn{K+1}{}
partition-specific coefficient vectors constrained to smoothness) and
non-interactive linear terms (``flat'' terms, specified via
\code{just\_linear\_without\_interactions}, receiving a single shared
coefficient vector \eqn{\mathbf{v}}{} across all partitions), the standard
solver carries \eqn{K+1}{} copies of \eqn{\mathbf{v}}{} linked by equality
constraints. Backfitting avoids this inflation by solving a
lower-dimensional problem at each step. Write the partition-\eqn{k}{}
design as
\eqn{\mathbf{X}_k = [\mathbf{Z}_k \mid
\mathbf{X}_{\mathrm{flat}}^{(k)}]}{}, where \eqn{\mathbf{Z}_k}{} contains
the spline columns and \eqn{\mathbf{X}_{\mathrm{flat}}^{(k)}}{} the flat
columns. This is invoked when \code{blockfit = TRUE}, flat columns are
non-empty, and \eqn{K > 0}{}.

The design, penalty, and constraint matrices are split into spline and
flat components by \code{\LinkA{.bf\_split\_components}{.bf.Rul.split.Rul.components}}, which extracts
spline rows from \eqn{\mathbf{A}}{}, drops null columns, rank-reduces via
QR, and detects mixed constraints (columns of \eqn{\mathbf{A}}{} with
nonzero entries on both spline and flat rows).
\end{SubSection}


%
\begin{SubSection}{Block-Coordinate Descent}
\strong{Spline step.} Holding \eqn{\mathbf{v}}{} fixed, the code forms
the adjusted response
\eqn{\mathbf{y}_k - \mathbf{X}_{\mathrm{flat}}^{(k)}\mathbf{v}}{} and
applies the spline-only Lagrangian projection via
\code{\LinkA{.bf\_lagrangian\_project}{.bf.Rul.lagrangian.Rul.project}}:
\deqn{\tilde{\boldsymbol{\beta}}_{\mathrm{spline}}^{(k)} =
  \mathbf{U}_{\mathrm{spline}}\mathbf{G}_{\mathrm{spline}}
  \mathbf{Z}_k^{\top}(\mathbf{y}_k -
  \mathbf{X}_{\mathrm{flat}}^{(k)}\mathbf{v}).}{}

\strong{Flat step.} Holding spline coefficients fixed, the shared flat
vector is updated by pooled penalized regression on residuals:
\deqn{\mathbf{v} = \left(\sum_{k=0}^{K}
  \mathbf{X}_{\mathrm{flat}}^{(k)\top}
  \mathbf{X}_{\mathrm{flat}}^{(k)} +
  \boldsymbol{\Lambda}_{\mathrm{flat}}\right)^{-1}
  \sum_{k=0}^{K}
  \mathbf{X}_{\mathrm{flat}}^{(k)\top}
  (\mathbf{y}_k - \mathbf{Z}_k
  \tilde{\boldsymbol{\beta}}_{\mathrm{spline}}^{(k)}).}{}
When the constraint matrix \eqn{\mathbf{A}}{} has mixed columns (nonzero
entries on both spline and flat rows), the flat update instead solves a
KKT system enforcing the residual equality constraint
\eqn{\mathbf{A}_{\mathrm{flat}}^{\top}\mathbf{v} = \mathbf{c} -
\mathbf{A}_{\mathrm{spline}}^{\top}
\tilde{\boldsymbol{\beta}}_{\mathrm{spline}}}{}
via \code{\LinkA{.bf\_constrained\_flat\_update}{.bf.Rul.constrained.Rul.flat.Rul.update}}.

Convergence is checked using the maximum absolute change across both
spline and flat coefficients. In the weighted inner loop used by the GLM
solvers, the flat-block change alone determines stopping.
\end{SubSection}


%
\begin{SubSection}{Four Estimation Cases}
\code{\LinkA{blockfit\_solve}{blockfit.Rul.solve}} dispatches to one of four paths.

\strong{Case (a): Gaussian identity + GEE}
(\code{\LinkA{.bf\_case\_gauss\_gee}{.bf.Rul.case.Rul.gauss.Rul.gee}}).
Whitening destroys block-diagonal structure, so the code skips
backfitting and performs the same full-system Gaussian GEE projection
used by \code{\LinkA{get\_B}{get.Rul.B}} Path 1a. The result is split back into
spline and flat components for downstream assembly.

\strong{Case (b): Gaussian identity, no correlation}
(\code{\LinkA{.bf\_case\_gauss\_no\_corr}{.bf.Rul.case.Rul.gauss.Rul.no.Rul.corr}}).
Standard block-coordinate descent as described above. The spline-only
\eqn{\mathbf{G}_{\mathrm{spline}}}{} factors are precomputed once, and the
pooled flat penalized inverse is reused across iterations.

\strong{Case (c): GLM + GEE}
(\code{\LinkA{.bf\_case\_glm\_gee}{.bf.Rul.case.Rul.glm.Rul.gee}}).
Two stages. Stage 1 forms a warm start by running damped Newton-Raphson
on the unwhitened working response: each outer iteration computes
working responses and weights at the current linear predictor, then
calls \code{\LinkA{.bf\_inner\_weighted}{.bf.Rul.inner.Rul.weighted}} for the inner backfitting loop.
Stage 2 refines the warm start on the full whitened system using the
damped SQP loop (\code{\LinkA{.bf\_sqp\_loop}{.bf.Rul.sqp.Rul.loop}}), replicating the approach
used by \code{\LinkA{.get\_B\_gee\_glm}{.get.Rul.B.Rul.gee.Rul.glm}}.

\strong{Case (d): GLM without GEE}
(\code{\LinkA{.bf\_case\_glm\_no\_corr}{.bf.Rul.case.Rul.glm.Rul.no.Rul.corr}}).
A damped Newton-Raphson outer loop updates working responses and weights,
while each inner iteration alternates between weighted spline and
weighted flat updates via \code{\LinkA{.bf\_inner\_weighted}{.bf.Rul.inner.Rul.weighted}}. Deviance is
monitored across outer iterations for convergence and damping.
\end{SubSection}


%
\begin{SubSection}{Inequality Constraints and Reassembly}
Because flat coefficients are shared by construction, the corresponding
equality constraints are satisfied exactly. Smoothness constraints on the
spline block are enforced by the spline-only Lagrangian projection.
After backfitting convergence, inequality constraints are enforced via
the same partition-wise active-set or dense SQP refinement used by
\code{\LinkA{get\_B}{get.Rul.B}}. The method is selected automatically by
\code{.solver\_detect\_qp\_global}: block-separable constraints use the
active-set method through \code{\LinkA{.active\_set\_refine}{.active.Rul.set.Rul.refine}}, while
cross-partition constraints trigger the dense SQP loop through
\code{\LinkA{.bf\_sqp\_loop}{.bf.Rul.sqp.Rul.loop}}. For GEE (Case c), inequality handling
occurs inside Stage 2 on the whitened system.

After convergence, the shared flat vector \eqn{\mathbf{v}}{} is copied
into each partition's coefficient vector, yielding
\eqn{\boldsymbol{\beta}_k =
[\tilde{\boldsymbol{\beta}}_{\mathrm{spline}}^{(k)\top},
\mathbf{v}^{\top}]^{\top}}{} for compatibility with downstream inference.
If \code{\LinkA{blockfit\_solve}{blockfit.Rul.solve}} throws an error, a warning is issued
and the code falls back to \code{\LinkA{get\_B}{get.Rul.B}}.
\end{SubSection}

\end{Section}
%
\begin{Section}{Knot Selection and Partitioning}


The topic of knot selection is not the main focus of the package, but the
partition structure is central because every later design matrix, penalty,
and smoothness constraint depends on it. The defaults in \pkg{lgspline} are
therefore meant to be practical and transparent rather than theoretically
final.

%
\begin{SubSection}{Univariate Case}
For a single predictor, the default partitioning is now handled by
\code{\LinkA{make\_partitions}{make.Rul.partitions}} in the same \eqn{k}{}-means framework used more
generally: \eqn{K+1}{} centers are fit on an internally standardized copy
of the predictor, controlled by \code{standardize\_predictors\_for\_knots}, and
then returned on the raw scale. Custom knots can still be supplied via
\code{custom\_knots}, in which case partition assignment is built directly
from those raw-scale breakpoints. The default number of knots \eqn{K}{} is
chosen adaptively based on \eqn{N}{}, \eqn{p}{}, \eqn{q}{}, and the GLM family.
For multivariate fits, the resulting partition metadata are returned as
\code{make\_partition\_list} and can be re-used in later calls to
\code{\LinkA{lgspline}{lgspline}}. This is particularly useful when one wants to hold
the partition geometry fixed across repeated fits, for example while varying
penalties, families, or correlation structures.
\end{SubSection}

%
\begin{SubSection}{Multivariate Case}
For multiple predictors, \eqn{K+1}{} cluster centers are identified by
\eqn{k}{}-means on an internally standardized predictor matrix via
\code{\LinkA{make\_partitions}{make.Rul.partitions}}. This is the partitioning mechanism used to
determine the multivariate spline regions; see MacQueen (1967) for the
classical clustering formulation and Kisi et al. (2025) for a recent applied
example of \eqn{k}{}-means-driven partitioning in a nonlinear prediction
setting. Midpoints between neighboring centers
(those whose midpoint  does not fall into a third cluster) serve as knot
locations. Observations are assigned to the nearest cluster center using
\code{\LinkA{get.knnx}{get.knnx}}, and the returned centers and knots are on
the original predictor scale. The resulting partition structure is a type
of Voronoi diagram and is stored in the fitted object as
\code{make\_partition\_list}. The \code{do\_not\_cluster\_on\_these} argument can
exclude certain predictors from clustering (e.g., a treatment indicator that
should not drive partitioning). The lower-level clustering behavior can be
further controlled by \code{cluster\_args} and \code{neighbor\_tolerance},
while \code{cluster\_on\_indicators} determines whether binary predictors are
allowed to influence the partition geometry at all.
\end{SubSection}


%
\begin{SubSection}{Standardizing Predictors}
Higher-order polynomial terms can dramatically inflate or deflate the
magnitude of basis expansions, introducing numerical instability. All
polynomial basis expansions are scaled by
\eqn{q_{0.69} - q_{0.31}}{}, where \eqn{q_{\zeta}}{} is the \eqn{\zeta}{}-th
quantile of the expansion. For a standard normal distribution this quantity
is approximately 1, so the scaling is close to one standard deviation for
symmetric distributions. This fitting-stage rescaling is controlled by
\code{standardize\_expansions\_for\_fitting}, while knot construction is
controlled separately by \code{standardize\_predictors\_for\_knots}. The same
scaling is applied to the constraint matrix to maintain smoothness, and
coefficients are back-transformed to the original scale after fitting.
\end{SubSection}

\end{Section}
%
\begin{Section}{Smoothing Spline Penalty}


%
\begin{SubSection}{Penalty Construction}
The penalty matrix \eqn{\boldsymbol{\Lambda}_s}{} penalizes the integrated
squared total curvature of the fitted function over the observed predictor
ranges. This is the step that makes the piecewise polynomial fit genuinely
behave like a smoothing spline rather than merely a constrained regression
spline. The package computes this penalty directly from the monomial
structure of the basis rather than by appealing to a pre-tabulated spline
basis. For a single partition \eqn{k}{} with basis expansion
\eqn{\mathbf{x}_k = (\phi_1(\mathbf{t}), \ldots, \phi_p(\mathbf{t}))^{\top}}{}
where each \eqn{\phi_i(\mathbf{t}) = \prod_{j=1}^{q} t_j^{\alpha_{ij}}}{}
is a multivariate monomial:
\deqn{\boldsymbol{\beta}_k^{\top}\boldsymbol{\Lambda}_s\boldsymbol{\beta}_k
  = \int_{\mathbf{a}}^{\mathbf{b}}
    \|\tilde{f}_k''(\mathbf{t})\|^{2}\,d\mathbf{t},}{}
where \eqn{\mathbf{a}}{} and \eqn{\mathbf{b}}{} are the observed predictor
minimums and maximums (computed globally from the data, not
partition-specific), and
\eqn{\tilde{f}_k(\mathbf{t}) = \mathbf{x}_k^{\top}\boldsymbol{\beta}_k}{}
is the fitted function for partition \eqn{\mathcal{P}_k}{}.

\strong{Total curvature operator.}
The integrated squared second derivative decomposes into \eqn{q}{}
curvature operators, one per predictor. For predictor \eqn{v}{}, the
curvature operator \eqn{D_v}{} is defined as
\deqn{D_v = \frac{\partial^{2}}{\partial t_v^{2}}
  + \sum_{s \neq v}\frac{\partial^{2}}{\partial t_v\,\partial t_s}.}{}
That is, \eqn{D_v}{} captures both the pure second derivative with respect
to \eqn{t_v}{} and all mixed second partial derivatives involving \eqn{t_v}{}.
The penalty matrix entries are then
\deqn{[\boldsymbol{\Lambda}_s]_{ij}
  = \sum_{v=1}^{q}\int_{\mathbf{a}}^{\mathbf{b}}
    D_v(\phi_i)\,D_v(\phi_j)\,d\mathbf{t}.}{}

\strong{Monomial derivative rule.}
For a monomial \eqn{\phi(\mathbf{t}) = \prod_j t_j^{\alpha_j}}{}, the
derivatives entering \eqn{D_v}{} have closed forms. The pure second
derivative is
\deqn{\frac{\partial^{2}}{\partial t_v^{2}}\prod_j t_j^{\alpha_j}
  = \alpha_v(\alpha_v - 1)\,t_v^{\alpha_v - 2}\prod_{j \neq v} t_j^{\alpha_j},}{}
which is zero when \eqn{\alpha_v < 2}{}. The mixed second derivative is
\deqn{\frac{\partial^{2}}{\partial t_v\,\partial t_s}\prod_j t_j^{\alpha_j}
  = \alpha_v\alpha_s\,t_v^{\alpha_v - 1}t_s^{\alpha_s - 1}
  \prod_{j \neq v,s} t_j^{\alpha_j},}{}
which is zero when \eqn{\alpha_v < 1}{} or \eqn{\alpha_s < 1}{}. Applying
\eqn{D_v}{} to a monomial \eqn{\phi_i}{} produces a sum of monomials with
known coefficients and exponent vectors.

\strong{Factorized integration.}
Because every \eqn{D_v(\phi_i)}{} is polynomial, the product
\eqn{D_v(\phi_i)\,D_v(\phi_j)}{} is also polynomial and the multivariate
integral factorizes over predictors:
\deqn{\int_{\mathbf{a}}^{\mathbf{b}}\prod_{j=1}^{q} t_j^{e_j}\,d\mathbf{t}
  = \prod_{j=1}^{q}\frac{b_j^{e_j+1} - a_j^{e_j+1}}{e_j + 1}.}{}
Crucially, this integral runs over \emph{all} \eqn{q}{} predictor ranges,
including predictors that do not appear in the integrand (for which
\eqn{e_j = 0}{} and the factor reduces to \eqn{b_j - a_j}{}). This ensures
that the penalty is properly scaled relative to the volume of the
predictor space.

\strong{Single-predictor verification.}
For \eqn{q = 1}{} with expansion
\eqn{\mathbf{x} = (1, t, t^{2}, t^{3})^{\top}}{} on \eqn{[a, b]}{}, the
curvature operator reduces to \eqn{D_1 = \partial^{2}/\partial t^{2}}{} (no
mixed partials exist), and the penalty matrix reduces to
\deqn{\boldsymbol{\Lambda}_s
  = \int_a^b \mathbf{x}''\mathbf{x}''^{\top}\,dt
  = \begin{pmatrix}
      0 & 0 & 0 & 0 \\
      0 & 0 & 0 & 0 \\
      0 & 0 & 4(b - a) & 6(b^{2} - a^{2}) \\
      0 & 0 & 6(b^{2} - a^{2}) & 12(b^{3} - a^{3})
    \end{pmatrix},}{}
equivalent to the classical cubic smoothing spline penalty originally proposed by Reinsch.

\strong{Handling of non-spline predictors.}
Predictors specified via \code{just\_linear\_without\_interactions} or
\code{just\_linear\_with\_interactions} do not receive higher-order
polynomial expansions in the design matrix. To ensure their curvature
contributions are still correctly computed (particularly through
interaction terms), the implementation temporarily appends phantom
higher-order columns (with zero data) for these predictors, computes
the full curvature penalty on the augmented basis, and then subsets the
result back to the original \eqn{p \times p}{} dimensions. This ensures
that interaction terms involving non-spline predictors receive appropriate
penalty contributions without affecting the rest of the estimation
pipeline.

\strong{Parallel computation.}
Because the total penalty is an additive sum over predictors
(\eqn{\boldsymbol{\Lambda}_s = \sum_{v=1}^{q}\boldsymbol{\Lambda}_{s,v}}{}),
the computation can be parallelized by distributing the per-predictor
curvature matrices across workers via \code{parallel::parLapply} and
summing the results. This is controlled by the \code{parallel\_penalty}
argument and is beneficial when \eqn{q}{} is large.

The penalty is computed by \code{\LinkA{get\_2ndDerivPenalty}{get.Rul.2ndDerivPenalty}} (single
predictor or subset) and \code{\LinkA{get\_2ndDerivPenalty\_wrapper}{get.Rul.2ndDerivPenalty.Rul.wrapper}}
(full assembly with optional parallelism and non-spline handling).

Because the smoothing penalty has zero eigenvalues for the intercept and
linear terms (whose second derivatives vanish), an optional ridge penalty on
lower-order terms is added for computational stability. The full penalty
block for partition \eqn{k}{} is:
\deqn{\boldsymbol{\Lambda}_k = \lambda_w\bigl(\boldsymbol{\Lambda}_s + \lambda_r\boldsymbol{\Lambda}_r + \sum_{m=1}^{M}\xi_{mk}\mathbf{L}_{mk}\bigr)}{}
where \eqn{\lambda_w}{} is the global wiggle penalty (\code{wiggle\_penalty}),
\eqn{\lambda_r}{} is ridge penalty on linear and intercept terms
(\code{flat\_ridge\_penalty}) multiplied by the wiggle penalty, and
\eqn{\xi_{mk}}{} and \eqn{\mathbf{L}_{mk}}{} denote optional additional
penalty multipliers and matrices, including the predictor- and
partition-specific components activated through
\code{unique\_penalty\_per\_predictor}, \code{unique\_penalty\_per\_partition},
\code{predictor\_penalties}, and \code{partition\_penalties}. This assembly is handled by
\code{\LinkA{compute\_Lambda}{compute.Rul.Lambda}}.

The penalty matrix \eqn{\boldsymbol{\Lambda}}{} is stored as a list
of \eqn{K+1}{} \eqn{p \times p}{} square, symmetric, positive semi-definite
matrices.
\end{SubSection}


%
\begin{SubSection}{Penalty Optimization via Leave-One-Out or Generalized Cross-Validation}
After the structural pieces of the model are fixed, the main remaining
question is how much smoothing to apply. In \pkg{lgspline}, that tuning is
performed with exact leave-one-out by default, or with generalized
cross-validation when \code{tuning\_criterion = "gcv"}. In either case, the
criterion is evaluated using the same constrained estimator that will be
used in the final model fit.
Penalty parameters are estimated on the log scale via exponential
parameterization (\eqn{\lambda = \exp(\theta)}{},
\eqn{\theta \in \mathbb{R}}{}), ensuring positivity. The chain rule factor
\eqn{\partial\exp(\theta)/\partial\theta = \exp(\theta) = \lambda}{} is
applied throughout. User-facing arguments (\code{initial\_wiggle},
\code{initial\_flat}, \code{predictor\_penalties},
\code{partition\_penalties}) accept values on the raw, natural scale;
conversion to log scale is handled internally. The final tuned values and
assembled penalty pieces are returned in the fitted object's
\code{penalties} component.

The total penalty matrix \eqn{\boldsymbol{\Lambda}}{} is constructed as:
\deqn{\boldsymbol{\Lambda} = \lambda_w\mathbf{L}_{w} + \lambda_r\mathbf{L}_{r}
  + \sum_{j}\nu_j\mathbf{L}_j^{(\mathrm{pred})}
  + \sum_{k}\tau_k\mathbf{L}_k^{(\mathrm{part})},}{}
where \eqn{\mathbf{L}_{w}}{} is the integrated squared second-derivative
penalty (i.e., \eqn{\boldsymbol{\Lambda}_s}{} above), \eqn{\mathbf{L}_{r}}{} is
a ridge penalty on intercept and linear coefficients,
\eqn{\mathbf{L}_j^{(\mathrm{pred})}}{} are predictor-specific penalties, and
\eqn{\mathbf{L}_k^{(\mathrm{part})}}{} are partition-specific penalties. The
scalars \eqn{\lambda_w}{} (\code{wiggle\_penalty}), \eqn{\lambda_r}{}
(\code{flat\_ridge\_penalty}), \eqn{\{\nu_j\}}{}
(\code{predictor\_penalties}), and \eqn{\{\tau_k\}}{}
(\code{partition\_penalties}) are tuned.
With \code{tuning\_criterion = "loo"}, the criterion is
\deqn{\mathrm{LOO}
  = \frac{1}{N}\sum_{i=1}^{N}\left(\frac{r_i}{1-h_{ii}}\right)^2,}{}
where \eqn{h_{ii}}{} is the diagonal of the tuning hat matrix. For Gaussian
identity-link tuning with fixed transformed design, this is exact. With
correlation structure and/or non-identity links, \pkg{lgspline} applies the
same transformed or linearized tuning problem already used by GCV, so the
leave-one-out calculation is carried out on that working scale. In empirical
diagnostics, the criterion itself behaves stably, while the fully analytic
derivative of the observation-wise leverage term can be numerically delicate
in some problems; when that sensitivity matters in practice, the finite-
difference optimizer path (\code{use\_custom\_bfgs = FALSE}) remains
available. This exact observation-wise calculation is also the main reason
LOO tuning becomes more expensive than GCV as \eqn{N}{} grows; in routine use,
\code{tuning\_criterion = "gcv"} is often the more practical choice once the
sample size is above about 250,000.

With \code{tuning\_criterion = "gcv"}, the package instead uses
\deqn{\mathrm{GCV}_{u,\gamma}
  = \frac{\sum_{i=1}^{N}D_{ii}\,r_i^{2}}
          {N(1 - \gamma\bar{W})^{2}}}{}
where \eqn{\bar{W} = \mathrm{tr}(\mathbf{H})/N}{} is the mean of the
hat-matrix diagonal. The symbol \eqn{\gamma}{} is set by
\code{gcv\_gamma} and controls optional inflation of the
effective-degrees-of-freedom term.

For identity link, \eqn{r_i = y_i - \hat{\eta}_i}{}. For non-identity links,
the residuals are
\eqn{r_i = g((y_i + \delta)/(1+2\delta)) - (\hat{\eta}_i + \delta)/(1+2\delta)}{},
where \eqn{\delta \geq 0}{} is a pseudocount that stabilizes the link
transformation, automatically tuned within \code{\LinkA{tune\_Lambda}{tune.Rul.Lambda}} if not supplied
to \code{delta}.
Non-Gaussian families with observation weights
\eqn{\omega_i}{} have their residuals scaled by \eqn{\omega_i}{}. When the
family provides a custom deviance residual function, that function is used
in place of the link-scale residuals.

\strong{Pseudocount selection.} The pseudocount \eqn{\delta}{} is chosen to
make the transformed response distribution most closely approximate a
\eqn{t}{}-distribution with \eqn{N-1}{} degrees of freedom, in the sense of
minimizing the (optionally weighted) mean absolute deviation between the
sorted standardized transformed responses and the corresponding
\eqn{t}{}-quantiles. This is solved via Brent's method over
\eqn{[10^{-64}, 1]}{}. When the link is identity, or when the response
naturally lies in the domain of the link function, \eqn{\delta = 0}{}.
This behavior is exposed through the \code{delta} argument in
\code{\LinkA{lgspline}{lgspline}}: supplying a fixed numeric value bypasses the internal
search, while leaving it \code{NULL} allows the tuning code to choose the
stabilizing pseudocount automatically when needed.

\strong{Meta-penalty regularization.} A regularization term pulls the
predictor- and partition-specific penalty parameters toward 1 on the raw
scale:
\deqn{P_{\mathrm{meta}}(\lambda_w, \nu_j, \tau_k)
  = \frac{1}{2}c_{\mathrm{meta}}\sum_j(\nu_j - 1)^{2}
  + \frac{1}{2}\cdot 10^{-32}(\lambda_w - 1)^{2}}{}
where \eqn{c_{\mathrm{meta}}}{} is a user-specified coefficient
(\code{meta\_penalty}). The gradient of \eqn{P_{\mathrm{meta}}}{} on the log
scale, incorporating the exp chain rule, is
\eqn{\partial P_{\mathrm{meta}}/\partial\theta_j = c_{\mathrm{meta}}(\nu_j - 1)\nu_j}{}
and
\eqn{\partial P_{\mathrm{meta}}/\partial\theta_1 = 10^{-32}(\lambda_w - 1)\lambda_w}{}.
The total objective is the selected tuning criterion plus
\eqn{P_{\mathrm{meta}}}{}.
\end{SubSection}


%
\begin{SubSection}{Closed-Form Gradients for Tuning Criteria}
The gradient of the GCV objective with respect to
\eqn{\theta_1 = \log\lambda_w}{} is computed analytically via the quotient
rule:
\deqn{\frac{\partial\mathrm{GCV}_{u,\gamma}}{\partial\theta_1}
  = \frac{1}{D^{2}}\left(\frac{\partial\mathcal{N}}{\partial\theta_1}D
  - \mathcal{N}\frac{\partial D}{\partial\theta_1}\right),}{}
where \eqn{\mathcal{N} = \sum r_i^{2}}{} (numerator) and
\eqn{D = N(1 - \gamma\bar{W})^{2}}{} (denominator). The key intermediates are:
\begin{itemize}

\item{} \eqn{\partial\mathbf{G}/\partial\lambda_w}{}, computed from the
matrix identity
\eqn{\partial(\mathbf{X}^{\top}\mathbf{X} + \boldsymbol{\Lambda})^{-1}/\partial\lambda = -\mathbf{G}(\partial\boldsymbol{\Lambda}/\partial\lambda)\mathbf{G}}{}.
\item{} \eqn{\partial\mathbf{G}^{1/2}/\partial\lambda_w}{}, derived from
\eqn{\partial\mathbf{G}/\partial\lambda_w}{} via the eigendecomposition
chain rule.
\item{} \eqn{\partial\bar{W}/\partial\lambda_w}{}, the derivative of the
trace of the hat matrix
\eqn{\mathbf{H} = \mathbf{X}\mathbf{U}\mathbf{G}\mathbf{X}^{\top}}{},
which depends on both
\eqn{\partial\mathbf{G}/\partial\lambda_w}{} and
\eqn{\partial\mathbf{G}^{1/2}/\partial\lambda_w}{}.
\item{} \eqn{\partial\mathcal{N}/\partial\theta_1 = -2\mathbf{r}^{\top}\mathbf{X}(\partial(\mathbf{U}\mathbf{G})/\partial\lambda_w)\mathbf{X}^{\top}\mathbf{y}\cdot\lambda_w}{},
via the chain rule applied to the residual vector.
\item{} \eqn{\partial D/\partial\theta_1 =
    2(1 - \gamma\bar{W})(-\gamma\,\partial\bar{W}/\partial\lambda_w)\cdot\lambda_w}{}.

\end{itemize}

In the implementation, these quantities are assembled by a small set of
helper routines: \code{\LinkA{compute\_dG\_dlambda}{compute.Rul.dG.Rul.dlambda}} for
\eqn{\partial\mathbf{G}/\partial\lambda}{}, \code{\LinkA{compute\_dGhalf}{compute.Rul.dGhalf}}
for \eqn{\partial\mathbf{G}^{1/2}/\partial\lambda}{},
\code{\LinkA{compute\_trace\_H}{compute.Rul.trace.Rul.H}} together with the tuning hat-diagonal
derivative helpers for the effective-degrees-of-freedom term, and
\code{\LinkA{compute\_dG\_u\_dlambda\_xy}{compute.Rul.dG.Rul.u.Rul.dlambda.Rul.xy}} for the derivative of the fitted-value
quadratic form.

For exact leave-one-out, the implementation instead differentiates the
per-observation PRESS form directly, requiring both
\eqn{\partial r_i / \partial\lambda_w}{} and
\eqn{\partial h_{ii} / \partial\lambda_w}{}. The latter is assembled
partition-wise from the same constrained-\eqn{\mathbf{G}}{} ingredients used
elsewhere in tuning. The exact LOO path does not explicitly form the full
projection matrix \eqn{\mathbf{U}}{} or the full hat matrix \eqn{\mathbf{H}}{};
instead it works with blockwise matrices
\eqn{\mathbf{C}_k = \mathbf{G}_k - \mathbf{G}_k\mathbf{A}_k(\mathbf{A}^\top\mathbf{G}\mathbf{A})^{-1}\mathbf{A}_k^\top\mathbf{G}_k}{}
and their derivatives. Empirically, this observation-wise leverage
derivative is the most numerically delicate part of the analytic LOO
gradient: most observations are typically well behaved, but a diffuse set of
small local errors can remain even when the overall LOO objective and the
corresponding GCV gradient are stable.

The shared wiggle and flat directions are differentiated directly. For
predictor- and partition-specific penalties, a lower-cost ratio heuristic is
used:
\deqn{\frac{\partial\mathrm{GCV}_{u,\gamma}}{\partial\theta_l} \approx
  \frac{\mathrm{mean}(\mathrm{diag}(\mathbf{L}_l))}{\mathrm{mean}(\mathrm{diag}(\boldsymbol{\Lambda}))}
  \frac{\partial\mathrm{GCV}_{u,\gamma}}{\partial\theta_w},}{}
where \eqn{\theta_l = \log \lambda_l}{} and
\eqn{\mathbf{L}_l}{} denotes the current contribution of the corresponding
predictor- or partition-specific penalty to \eqn{\boldsymbol{\Lambda}}{}. This
keeps the extra penalty directions inexpensive while preserving the direct
derivatives for the shared wiggle and flat penalties.
\end{SubSection}


%
\begin{SubSection}{Optimization Procedure}
\strong{Grid search initialization.} The selected tuning criterion is
evaluated over a grid of candidate values for
\eqn{(\lambda_w, \lambda_r)}{} on the log scale. All combinations of
user-supplied candidate vectors (\code{initial\_wiggle} and
\code{initial\_flat}) are formed, and the combination yielding
the smallest finite criterion value is selected as the starting
point for BFGS optimization. Grid points producing non-finite
objective values are discarded. If all grid points fail, an error
is raised advising the user to check the data or adjust the grid.

\strong{Damped BFGS optimizer.} A custom damped BFGS quasi-Newton optimizer,
implemented in \code{\LinkA{tune\_Lambda}{tune.Rul.Lambda}} through
\code{.damped\_bfgs()}, minimizes the selected tuning criterion plus
\eqn{P_{\mathrm{meta}}}{}. When analytic gradients are not preferred, or when
the exact LOO leverage derivative appears too noisy for a given problem, a
base-R \code{\LinkA{optim}{optim}} call with method \code{"BFGS"} provides
the finite-difference fallback.

\emph{Iterations 1-2: steepest descent.} The first two iterations use
steepest descent with a damping factor \eqn{\alpha}{}:
\eqn{\boldsymbol{\phi}^{(t+1)} = \boldsymbol{\phi}^{(t)} - \alpha\nabla_{\boldsymbol{\phi}}}{}.

\emph{Iterations 3+: BFGS.} From iteration 3, an inverse Hessian
approximation \eqn{\mathbf{J}^{(t)}}{} is maintained via the standard secant
update. Let
\eqn{\mathbf{s}^{(t)} = \boldsymbol{\phi}^{(t)} - \boldsymbol{\phi}^{(t-1)}}{}
and
\eqn{\mathbf{v}^{(t)} = \nabla^{(t)} - \nabla^{(t-1)}}{}. The BFGS update
is:
\deqn{\mathbf{J}^{(t+1)}
  = (\mathbf{I} - \mathbf{u}\mathbf{s}\mathbf{v}^{\top})
  \mathbf{J}^{(t)}
  (\mathbf{I} - \mathbf{u}\mathbf{v}\mathbf{s}^{\top})
  + \mathbf{u}\mathbf{s}\mathbf{s}^{\top},
  \qquad \mathbf{u} = (\mathbf{v}^{\top}\mathbf{s})^{-1}.}{}
When \eqn{|\mathbf{v}^{\top}\mathbf{s}| < 10^{-64}}{}, the approximation is
reset to \eqn{\mathbf{I}}{} and the iteration is flagged for restart. The
search direction is
\eqn{\mathbf{d}^{(t)} = -\mathbf{J}^{(t)}\nabla^{(t)}}{}.

\emph{Step acceptance.} A step is accepted if the new tuning criterion value
is no larger than the old one.
On rejection, \eqn{\alpha}{} is halved. If \eqn{\alpha < 2^{-10}}{} (early
iterations) or \eqn{\alpha < 2^{-12}}{} (later iterations), the optimizer
terminates with the best solution found.

\emph{Convergence.} The optimizer terminates when
the absolute change in the tuning criterion is less than \eqn{\epsilon}{}, or
\eqn{\|\boldsymbol{\phi}^{(t)} - \boldsymbol{\phi}^{(t-1)}\|_{\infty} < \epsilon}{},
provided at least 10 iterations have elapsed, for penalties
\eqn{\boldsymbol{\phi} = \lambda_w, \lambda_r, \nu_j, \tau_j, ...}{}.

\emph{Alternative.} A base-R
\code{\LinkA{optim}{optim}} call with method \code{"BFGS"} and
finite-difference gradients can be used instead via
\code{use\_custom\_bfgs = FALSE}.

\strong{Post-optimization sample-size adjustment.} After optimization, the
tuned penalty parameters are divided by \eqn{((N+2)/(N-2))^{2}}{}
(equivalently multiplied by \eqn{((N-2)/(N+2))^{2}}{}) for both GCV- and
LOO-based tuning. This decreases the final penalties at small sample sizes.

The tuning loop is implemented in \code{\LinkA{tune\_Lambda}{tune.Rul.Lambda}}.
\end{SubSection}

\end{Section}
%
\begin{Section}{Incorporating Non-Spline Effects}


Multiple fixed effects are accommodated naturally in the LMSS framework
because spline effects, linear effects, and many interaction terms all live
in the same partition-wise polynomial expansion. The distinction is therefore
not whether a term is "allowed" by the solver, but whether it receives full
spline treatment or remains structurally linear across partitions.

The constrained framework naturally accommodates non-spline terms. If only
linear terms are included for a predictor (via
\code{just\_linear\_without\_interactions} or
\code{just\_linear\_with\_interactions}), the first-derivative smoothing
constraint forces the linear coefficient to be identical across all
partitions, since the derivative of a linear function is its slope. This
is not an algorithmic modification but a natural consequence of the
constraint structure.

For example, a model with one spline effect and a linear treatment indicator
interaction will naturally keep the treatment-time interaction coefficient
constant across partitions while allowing the time effect to vary
nonlinearly. This conveniently extends to arbitrary combinations of spline
and linear terms without requiring special handling.

When \code{blockfit = TRUE} is specified alongside
\code{just\_linear\_without\_interactions}, the flat-block path provides an
alternative enforcement mechanism. Rather than relying on constraint
projection, flat coefficients are pooled structurally across partitions
during backfitting. The two approaches agree at the point estimate but
differ in their uncertainty quantification; see the Blockfit section above.
\end{Section}
%
\begin{Section}{Integration}


Because the fitted object retains an explicit polynomial representation in
each partition, numerical integration can be carried out in a fairly direct
way. The package wraps that calculation in a user-facing S3 method so the
user does not need to manage knot boundaries or partition membership by hand.

In the user-facing interface, numerical integration is applied through
\code{\LinkA{integrate.lgspline}{integrate.lgspline}}, which applies Gauss-Legendre quadrature to
predictions from the fitted model produced by \code{\LinkA{predict.lgspline}{predict.lgspline}}.

%
\begin{SubSection}{Implementation}
For a user-supplied rectangular domain, \code{\LinkA{integrate.lgspline}{integrate.lgspline}} constructs a
tensor-product grid of Gauss-Legendre nodes, evaluates the fitted model at
those points, and forms the weighted sum. This works for both univariate and
multivariate models, respects the fitted partition structure automatically,
and avoids requiring the user to keep track of knot boundaries by hand.

The \code{vars} argument selects which predictors are integrated over.
Predictors not listed in \code{vars} are held fixed at
\code{initial\_values} when supplied, or otherwise at the midpoint of their
observed training range. The optional \code{B\_predict} argument makes it
possible to integrate posterior draws or other alternate coefficient sets,
and \code{n\_quad} controls the number of Gauss-Legendre nodes used per
integrated dimension.

Integration is performed on the response scale by default. Setting
\code{link\_scale = TRUE} instead integrates the linear predictor
\eqn{\eta = f(\mathbf{t})}{}. For identity-link Gaussian models the two scales
coincide.
\end{SubSection}

\end{Section}
%
\begin{Section}{Lagrange Multipliers}


When \code{return\_lagrange\_multipliers = TRUE}, the multiplier vector
\deqn{\boldsymbol{\lambda} = (\mathbf{A}^{\top}\mathbf{G}\mathbf{A})^{-1}\mathbf{A}^{\top}\hat{\boldsymbol{\beta}}}{}
is returned. These quantify the sensitivity of the penalized objective to
relaxing each smoothness or user-supplied equality constraint. When
constraint target values are nonzero
(\eqn{\mathbf{A}^{\top}\boldsymbol{\beta}_0 \neq \mathbf{0}}{}), the modified
formulation is used:
\deqn{\boldsymbol{\lambda} = (\mathbf{A}^{\top}\mathbf{G}\mathbf{A})^{-1}\mathbf{A}^{\top}(\hat{\boldsymbol{\beta}} - \boldsymbol{\beta}_0)}{}
where \eqn{\mathbf{A}^{\top}\boldsymbol{\beta}_0}{} is the vector of
constraint target values. Multipliers are \code{NULL} when no constraints
are active (\eqn{\mathbf{A}}{} is \code{NULL} or \eqn{K = 0}{}).

For inequality constraints, multipliers are returned as computed by
\code{\LinkA{solve.QP}{solve.QP}}. The Lagrange multipliers for active
inequality constraints can be used diagnostically to identify which shape
constraints are most costly in terms of goodness of fit.
\end{Section}
%
\begin{Section}{S3 Methods}


Standard S3 methods are provided for objects of class \code{lgspline}:
\begin{itemize}

\item{} \code{\LinkA{print.lgspline}{print.lgspline}} and \code{\LinkA{summary.lgspline}{summary.lgspline}}:
Provide concise model summaries, with
\code{\LinkA{print.summary.lgspline}{print.summary.lgspline}} formatting coefficient tables in a
familiar regression-style layout.

\item{} \code{\LinkA{logLik.lgspline}{logLik.lgspline}}: Returns a standard \code{logLik}
object. For Gaussian responses with identity link, the exact
log-likelihood is computed. When a correlation structure is present via
\code{VhalfInv}, the log-likelihood includes the
\eqn{\log|\mathbf{V}^{-1/2}|}{} adjustment and the corresponding
whitened quadratic form. For other families, the method falls back to
\code{family\$aic} or a deviance-based approximation. An
\code{include\_prior} argument (default \code{TRUE}) optionally adds
the Gaussian prior penalty interpretation of the smoothing spline
penalty
\eqn{-\frac{1}{2\sigma^{2}}\tilde{\boldsymbol{\beta}}^{\top}\boldsymbol{\Lambda}\tilde{\boldsymbol{\beta}}}{}
to obtain a penalized MAP log-likelihood.

\item{} \code{\LinkA{predict.lgspline}{predict.lgspline}}: Produces fitted values and related
quantities (e.g., derivatives and standard errors through
\code{se.fit = TRUE}), lets \code{new\_predictors} override
\code{newdata}, accepts alternate coefficient lists through
\code{B\_predict}, and supports prediction on new predictor matrices
consistent with the original spline expansions.

\item{} \code{\LinkA{coef.lgspline}{coef.lgspline}}: Extracts partition-specific
coefficient vectors.

\item{} \code{\LinkA{confint.lgspline}{confint.lgspline}}: Extracts confidence intervals.
When the inverse Hessian from BFGS optimization is available for
correlation parameters, intervals for those correlation parameters are
returned on the working (transformed) scale and should be
back-transformed as described in the correlation section.
\item{} \code{\LinkA{plot.lgspline}{plot.lgspline}}: For one-dimensional fits, produces
base R graphics showing the fitted function (with optional partition-wise
formulas) and supports overlay via \code{add = TRUE}. For two or more
predictors, an interactive \pkg{plotly}-based visualization is
returned. Specific predictors may be selected via \code{vars}.

\item{} \code{\LinkA{integrate.lgspline}{integrate.lgspline}}: Computes definite integrals of
the fitted surface over rectangular domains by Gauss-Legendre
quadrature.

\end{itemize}


Additional user-facing helpers include \code{\LinkA{wald\_univariate}{wald.Rul.univariate}} for
coefficient-wise Wald inference, \code{\LinkA{generate\_posterior}{generate.Rul.posterior}} for
posterior and posterior-predictive sampling,
\code{\LinkA{generate\_posterior\_correlation}{generate.Rul.posterior.Rul.correlation}} for correlation-aware posterior
simulation, \code{\LinkA{equation}{equation}} for closed-form display of the fitted
partition formulas, and
\code{\LinkA{find\_extremum}{find.Rul.extremum}} for optimizing the fitted surface or a custom
acquisition function built from it.
\end{Section}
%
\begin{References}
Buse, A. and Lim, L. (1977). Cubic Splines as a Special Case of
Restricted Least Squares. \emph{Journal of the American Statistical
Association}, 72, 64-68.

Eilers, P. H. and Marx, B. D. (1996). Flexible Smoothing with B-splines
and Penalties. \emph{Statistical Science}, 11(2), 89-121.

Ezhov, N., Neitzel, F. and Petrovic, S. (2018). Spline Approximation,
Part 1: Basic Methodology. \emph{Journal of Applied Geodesy}, 12(2),
139-155.

Goldfarb, D. and Idnani, A. (1983). A Numerically Stable Dual Method for
Solving Strictly Convex Quadratic Programs. \emph{Mathematical
Programming}, 27(1), 1-33.

Harville, D. A. (1977). Maximum Likelihood Approaches to Variance
Component Estimation and to Related Problems. \emph{Journal of the
American Statistical Association}, 72(358), 320-338.

Hastie, T. J. and Tibshirani, R. J. (1990). \emph{Generalized Additive
Models}. Chapman \& Hall/CRC.

Kisi, O., Heddam, S., Parmar, K. S., Petroselli, A., Kulls, C. and
Zounemat-Kermani, M. (2025). Integration of Gaussian Process Regression
and K Means Clustering for Enhanced Short Term Rainfall Runoff Modeling.
\emph{Scientific Reports}, 15, 7444.

MacQueen, J. B. (1967). Some Methods for Classification and Analysis of
Multivariate Observations. In \emph{Proceedings of the Fifth Berkeley
Symposium on Mathematical Statistics and Probability}, Volume 1,
281-297. University of California Press.

McCullagh, P. and Nelder, J. A. (1989). \emph{Generalized Linear Models}.
Chapman \& Hall, 2nd edition.

Murray, I., Adams, R. P. and MacKay, D. J. C. (2010). Elliptical Slice
Sampling. \emph{Proceedings of the 13th International Conference on
Artificial Intelligence and Statistics (AISTATS)}, 9, 541-548.

Nocedal, J. and Wright, S. J. (2006). \emph{Numerical Optimization}
(2nd ed.). Springer.

Kim, Y.-J. and Gu, C. (2004). Smoothing Spline Gaussian Regression:
More Scalable Computation via Efficient Approximation. \emph{Journal of
the Royal Statistical Society: Series B}, 66(2), 337-356.

Patterson, H. D. and Thompson, R. (1971). Recovery of Inter-Block
Information When Block Sizes Are Unequal. \emph{Biometrika}, 58,
545-554.

Pya, N. and Wood, S. N. (2015). Shape Constrained Additive Models.
\emph{Statistics and Computing}, 25(3), 543-559.

Reinsch, C. H. (1967). Smoothing by Spline Functions. \emph{Numerische
Mathematik}, 10, 177-183.

Ruppert, D., Wand, M. P. and Carroll, R. J. (2003).
\emph{Semiparametric Regression}. Cambridge University Press.

Searle, S. R., Casella, G. and McCulloch, C. E. (2006). \emph{Variance
Components}. Wiley.

Wahba, G. (1990). \emph{Spline Models for Observational Data}. SIAM.

Wood, S. N. (2006). On Confidence Intervals for Generalized Additive
Models Based on Penalized Regression Splines. \emph{Australian \& New
Zealand Journal of Statistics}, 48(4), 445-464.

Wood, S. N. (2011). Fast Stable Restricted Maximum Likelihood and Marginal
Likelihood Estimation of Semiparametric Generalized Linear Models.
\emph{Journal of the Royal Statistical Society: Series B}, 73(1), 3-36.

Wood, S. N. (2017). \emph{Generalized Additive Models: An Introduction
with R}. CRC Press, 2nd edition.
\end{References}
